\documentclass[11pt,a4paper]{article}

\usepackage[utf8]{inputenc}
\usepackage[T1]{fontenc}
\usepackage{lmodern}
\usepackage[a4paper,margin=2.45cm]{geometry}
\usepackage{amsmath,amssymb,amsthm,mathtools}
\usepackage{booktabs}
\usepackage{array}
\usepackage{microtype}
\usepackage{enumitem}
\usepackage{xcolor}
\usepackage[hidelinks]{hyperref}
\usepackage{longtable}
\usepackage{graphicx}

\hypersetup{
  pdftitle={Causal Phase Obstructions in Exact Information Simulation},
  pdfauthor={Jan Mikulik}
}

\setcounter{MaxMatrixCols}{20}
\setlength{\parindent}{1.1em}
\setlength{\parskip}{0.18em}
\emergencystretch=2em
\setlist[itemize]{topsep=3pt,itemsep=2pt,parsep=1pt}
\setlist[enumerate]{topsep=3pt,itemsep=2pt,parsep=1pt}

\newtheorem{theorem}{Theorem}[section]
\newtheorem{proposition}[theorem]{Proposition}
\newtheorem{lemma}[theorem]{Lemma}
\newtheorem{corollary}[theorem]{Corollary}
\newtheorem{conjecture}[theorem]{Conjecture}
\newtheorem{problem}[theorem]{Open Problem}
\theoremstyle{definition}
\newtheorem{definition}[theorem]{Definition}
\theoremstyle{remark}
\newtheorem{remark}[theorem]{Remark}

\newcommand{\E}{\mathbb E}
\newcommand{\Pp}{\mathbb P}
\newcommand{\R}{\mathbb R}
\newcommand{\Z}{\mathbb Z}
\newcommand{\cP}{\mathcal P}
\newcommand{\cM}{\mathcal M}
\newcommand{\cK}{\mathcal K}
\newcommand{\cB}{\mathcal B}
\newcommand{\cD}{\mathfrak D}
\newcommand{\cR}{\mathcal R}
\newcommand{\cS}{\mathcal S}
\newcommand{\cW}{\mathcal W}
\newcommand{\id}{\operatorname{id}}
\newcommand{\supp}{\operatorname{supp}}
\newcommand{\Law}{\operatorname{Law}}
\newcommand{\Span}{\operatorname{span}}
\newcommand{\TV}{\operatorname{TV}}
\newcommand{\one}{\mathbf 1}
\newcommand{\bits}{\{0,1\}}
\newcommand{\wt}{\operatorname{wt}}
\newcommand{\diag}{\operatorname{diag}}

\title{\textbf{Causal Phase Obstructions in Exact Information Simulation}\\[0.35em]
\large Zero-Repair Trees, Hidden Ordering Information, and Multiscale Renormalization}
\author{Jan Mikul\'ik}
\date{23 August 2026\\[0.2em]\small Revised phase-fiber version 3}

\begin{document}
\maketitle

\begin{abstract}
We study an exact causal information-accounting problem in which two finite
information-increment laws must be reconciled by a nonnegative,
history-dependent inventory, with no fresh repair term.  The local constraint is
a linear convolution identity, but the global feasibility problem is strongly
order-sensitive.

Our first result is a finite-horizon converse.  For every fixed Laplace tilt we
derive an explicit logarithmic lower bound for the minimum root inventory
\(\cB_n\).  If the increment laws have equal mean and variance and their first
mismatch is an oriented third cumulant, then
\[
 \liminf_{n\to\infty}\frac{\cB_n}{\log n}
 \ge \frac{\Delta\kappa_3}{6\sigma^2}.
\]
For the dyadic pair
\[
 \mu_P=\tfrac14\delta_2+\tfrac34\delta_4,
 \qquad
 \mu_Q=\tfrac34\delta_3+\tfrac14\delta_5,
\]
this gives the exact lower-bound coefficient \(1/3\) and rules out every infinite
finite-mean zero-repair tree.

We then isolate a second obstruction which is invisible to scalar information
spectra.  At block length \(m\), ordered binary histories form a space
\(V_m\cong\R^{2^m}\), while scalar total information sees only the \(m+1\)
Hamming layers.  The lost ordering information is the phase space
\(W_m=\ker\Pi_m\), of dimension
\[
 \dim W_m=2^m-(m+1).
\]
Under block doubling it has the exact decomposition
\[
W_{2m}=(W_m\otimes W_m)\oplus(W_m\otimes U_m)
\oplus(U_m\otimes W_m)\oplus K_m,
\qquad \dim K_m=m^2,
\]
where \(K_m\) is a fresh anti-diagonal phase sector created by coarse-graining.
For two symbols we go further: the entire projected causal scalar-profile
polytope is the static transport polytope cut by exactly two additional linear
facets,
\[
2r_{00}+3r_{01}+r_{11}\ge\frac3{16},\qquad
2r_{00}+3r_{01}+2r_{10}+3r_{11}\le\frac98.
\]
This intrinsic criterion yields a solver-free proof that the minimum
information-profile mismatch is \(5/16\), whereas static matching pays only
\(1/4\); a fully displayed rational primal--dual certificate gives an
independent exact check.  More generally, at every fixed block length the
local causal phase fiber is a bounded polytope.  We give an exact
necessary-and-sufficient Farkas characterization of its nonemptiness; for
rational source laws the certificate cone is rational polyhedral and hence has
a finite rational generating family.  The extra \(1/16\) in the two-symbol
example is added to both scalar residual profiles and therefore disappears from
their signed difference.  Causality can thus create a real obstruction living
entirely in the kernel of scalar information projection.

Finally, we pass from viability sets to lower support functionals, derive an
exact finite-dimensional Bellman--Fenchel recursion, identify the cubic shift
normal form
\[
P^*-Q^*=\tfrac14E^2(I-E)^3,
\]
and introduce a Green coordinate in which the cubic defect becomes a positive
unit source.  The two-block boundary mask is \((1+z)^3/4\), the binary
quadratic B-spline mask, and one explicit mean-neutral phase channel admits a
uniform finite-mean envelope across all dyadic scales.  These statements do
not prove a matching \(O(\log n)\) upper bound.  Fixed-block phase feasibility
is now completely characterized at the level of finite-dimensional linear
programming; what remains is the genuinely multiscale problem of finding a
scale-stable description compatible with future zero-repair viability.  In
particular, a uniform bound on \(\cB_{2n}-\cB_n\) would close the order of
growth.

The zero-repair bill is deliberately distinguished from unrestricted causal
Shannon entropy.  It measures exact causal working inventory, not entropy rate;
this distinction is the link back to the broader question of which resources
an exact stochastic realization must carry.
\end{abstract}

\noindent\textbf{Keywords:} causal information, exact simulation, zero repair,
causal transport, information spectrum, renormalization, Bellman--Fenchel
duality, subdivision schemes, third cumulant.

\section{Introduction}

A probability law tells us what may be observed.  It does not, by itself, tell
us how much hidden state an exact realization must carry when outputs are
chosen sequentially and future repair is forbidden.  This paper isolates one
such resource question.

Suppose two actions have prescribed information-increment laws \(\mu_P\) and
\(\mu_Q\).  At every history we ask for an exact accounting identity: the
current nonnegative inventory plus the next \(Q\)-increment must have the same
law as a \(P\)-branch increment plus the next inventory assigned to that branch.
No signed debt is allowed and no fresh repair term may be injected.  The
resulting object is an exact zero-repair information tree.

There are two reasons to study this stricter contract separately from ordinary
entropy-rate simulation.  First, exact online random generation can recycle
unused randomness and operate arbitrarily close to Shannon entropy in models
that allow a persistent state and fresh source bits
\cite{DraperSaad2026,DraperHarrisSaad2026}.  Second, finitary coding theory
shows that equal entropy alone does not control stronger moments of coding
radius or information cocycles \cite{Parry1979,Schmidt1984,HarveyPeres2011}.
Thus rate and working state are different resources.

The present problem also sits near, but not inside for free, modern adapted
optimal transport.  Causal and bicausal transport admit general duality and
dual attainment under broad assumptions \cite{KrsekPammer2025}; graph-causal
transport now provides gluing and dynamic-programming results for couplings
constrained by directed acyclic graphs \cite{OblojTuchilus2026}.  Our model has
an additional exact conservation law and a hard state constraint
\(J\ge0\).  These are not cosmetic: they are the source of the obstruction.

A useful recent methodological parallel comes from synchronization channels.
Kova\v{c}evi\'c identifies a coefficientwise domination condition for a family
of sticky channels, converts it into an exact dual capacity bound, and then
characterizes the laws satisfying that condition through a functional
composition equation \cite{Kovacevic2026}.  The present setting is different,
but the same standard is appropriate: the objective is not another successful
ansatz, but an exact structural criterion exposing the admissible positive
fiber on which the dynamics acts.

\subsection{Main contributions}

The paper has six main components.

\begin{enumerate}
\item We give a self-contained finite-horizon logarithmic converse for exact
zero repair.  The third-cumulant case yields an explicit lower-bound coefficient, equal to
\(1/3\) for the dyadic parity pair.

\item We represent finite viability sets by their lower support functionals and
derive an exact Bellman--Fenchel recursion in finite dimensions, together with
a rigorous weak-duality version on the countable state space.

\item We introduce the ordered-history phase quotient.  Scalar total
information is a projection \(\Pi_m\) from \(2^m\) ordered histories to
\(m+1\) Hamming layers; its kernel has dimension \(2^m-(m+1)\).

\item We characterize the first nontrivial projected causal phase polytope in
closed form.  For two symbols a static scalar joint profile admits a causal
lift if and only if two explicit linear facet inequalities hold.  This gives a
solver-free proof of the exact mismatch \(5/16\), compared with the static
value \(1/4\), while an independent rational primal--dual certificate checks
the same optimum.  We then characterize the local causal phase fiber at
arbitrary fixed block length as a polytope and give an exact Farkas dual
criterion for its nonemptiness.

\item We prove a block-doubling decomposition of the phase kernel.  Besides
inherited phase directions, each doubling creates a fresh \(m^2\)-dimensional
anti-diagonal sector.  This rules out any globally one-dimensional phase
picture and shows why a fixed-block Farkas description is not yet a
scale-uniform renormalization theorem.

\item We identify the cubic finite-difference normal form, a Green-coordinate
affine recursion, the associated quadratic B-spline boundary mask, and one
explicit marginal phase channel that can be transported through all dyadic
scales at bounded mean cost.  The sharp upper bound remains open; the paper
ends with a precise renormalization target rather than a claimed QED.
\end{enumerate}

\subsection{Resource interpretation and scope}

The quantity studied here is not automatically the excess Shannon entropy of
an unrestricted causal simulator.  It is an auxiliary working-inventory cost
under a stricter contract.  We will return to this distinction in
Section~\ref{sec:original}.  In particular, no ontological conclusion is drawn
from the mathematics: the paper prices a stated realization constraint.

\section{Exact zero-repair information trees}
\label{sec:model}

Let
\[
 \mu_P=\sum_{i=1}^r p_i\delta_{x_i},
 \qquad p_i>0,
 \qquad \sum_i p_i=1,
\]
and let \(\mu_Q\) be a finite-support probability law on \([0,\infty)\).
Branch labels are retained even if two values \(x_i\) coincide; equal current
information does not imply equal future subtrees.

\begin{definition}[Depth-\(n\) exact zero-repair tree]
A depth-\(n\) zero-repair tree is a family of probability laws
\[
 \{B_h\}_{|h|\le n}
\]
on \([0,\infty)\), indexed by histories
\(h\in\{1,\ldots,r\}^{\le n}\), such that each internal node satisfies
\begin{equation}
 \boxed{
 \mu_Q*B_h
 =
 \sum_{i=1}^r p_i\,\delta_{x_i}*B_{hi}.
 }
 \label{eq:ZR}
\end{equation}
\end{definition}

For a buffer law \(B_h\), put \(m_h:=\E_{B_h}J\).

\begin{definition}[Finite-horizon zero-repair bill]
\begin{equation}
 \boxed{
 \cB_n:=\inf\left\{m_\varnothing:
 \{B_h\}_{|h|\le n}\text{ satisfies \eqref{eq:ZR}}\right\}.
 }
 \label{eq:bill}
\end{equation}
The value is \(+\infty\) if no feasible tree exists.
\end{definition}

Throughout the converse we assume equal increment means,
\begin{equation}
 \E_{\mu_P}X=\E_{\mu_Q}Y.
 \label{eq:equalmean}
\end{equation}

\begin{lemma}[Mean martingale]
If \(m_h<\infty\), then all child means are finite and
\begin{equation}
 \boxed{m_h=\sum_i p_i m_{hi}.}
 \label{eq:meanmart}
\end{equation}
Consequently, along an independently \(P\)-distributed branch history,
\(m_{H_t}\) is a nonnegative martingale.
\end{lemma}

\begin{proof}
Take first moments in \eqref{eq:ZR}.  Finite support and nonnegative child
means imply finiteness of each child mean.  The common increment mean cancels,
leaving \eqref{eq:meanmart}.
\end{proof}

\begin{remark}[Constrained coboundary form]
After coupling the equal laws in \eqref{eq:ZR}, one may write schematically
\[
 J_h+Y=x_I+J_{hI}.
\]
Thus \(Y-x_I=J_{hI}-J_h\) resembles a coboundary.  The difficulty is the
simultaneous imposition of nonnegativity, exact branch weights, convolutional
input, and future viability at every history.
\end{remark}

\section{A finite-horizon logarithmic converse}
\label{sec:lower}

Fix \(a>0\) and define
\[
 G_P(a):=\E_{\mu_P}e^{-aX},
 \qquad
 G_Q(a):=\E_{\mu_Q}e^{-aY},
 \qquad
 \rho(a):=\frac{G_Q(a)}{G_P(a)}.
\]
Introduce the tilted branch law
\[
 \widetilde p_i^{(a)}:=\frac{p_ie^{-ax_i}}{G_P(a)}
\]
and its likelihood-ratio second moment
\begin{equation}
 \boxed{
 \Lambda(a):=
 \sum_i\frac{(\widetilde p_i^{(a)})^2}{p_i}
 =\frac{G_P(2a)}{G_P(a)^2}\ge1.
 }
 \label{eq:Lambda}
\end{equation}
For each node set
\[
 \widehat B_h(a):=\E_{B_h}e^{-aJ},
 \qquad
 y_h:=e^{-am_h}.
\]
By Jensen and \(J\ge0\),
\begin{equation}
 0<y_h\le\widehat B_h(a)\le1.
 \label{eq:JensenBasic}
\end{equation}

For a \(k\)-step descendant \(hw\), define
\[
 D_h^{(k)}:=\E_{P^k}y_{hw}-y_h.
\]

\begin{lemma}[Jensen--Hellinger control]
For every node and every admissible \(k\),
\[
 D_h^{(k)}\ge0,
\]
and
\begin{equation}
 \E_{P^k}\left(\sqrt{y_{hw}}-\sqrt{y_h}\right)^2
 \le D_h^{(k)},
 \qquad
 \E_{P^k}(y_{hw}-y_h)^2\le4D_h^{(k)}.
 \label{eq:JH}
\end{equation}
\end{lemma}

\begin{proof}
By \eqref{eq:meanmart}, \(\E_{P^k}m_{hw}=m_h\).  Convexity of
\(u\mapsto e^{-au}\) gives \(D_h^{(k)}\ge0\), while
\[
 \E e^{-am_{hw}/2}\ge e^{-a\E m_{hw}/2}=\sqrt{y_h}.
\]
Expanding the square proves the first estimate.  Since
\(|u-v|\le2|\sqrt u-\sqrt v|\) on \([0,1]\), the second follows.
\end{proof}

Taking Laplace transforms in \eqref{eq:ZR} yields the exact tilted identity
\begin{equation}
 \boxed{
 \E_{\widetilde P_a}\widehat B_{hi}(a)
 =\rho(a)\widehat B_h(a),
 }
 \qquad
 \boxed{
 \E_{\widetilde P_a^k}\widehat B_{hw}(a)
 =\rho(a)^k\widehat B_h(a).
 }
 \label{eq:tiltedidentity}
\end{equation}
Let
\[
 L_w:=\frac{\widetilde P_a^k(w)}{P^k(w)}.
\]
Then \(\E L_w=1\) and \(\E L_w^2=\Lambda(a)^k\).

\begin{lemma}[One-block measure change]
For every node \(h\),
\begin{equation}
 \boxed{
 \rho(a)^k\widehat B_h(a)
 \ge y_h-2\Lambda(a)^{k/2}\sqrt{D_h^{(k)}}.
 }
 \label{eq:blocklocal}
\end{equation}
\end{lemma}

\begin{proof}
Since \(\E(L_w-1)=0\),
\[
 \E_{\widetilde P_a^k}y_{hw}
 =y_h+\E_{P^k}[(L_w-1)(y_{hw}-y_h)].
\]
Cauchy--Schwarz and \eqref{eq:JH} give the stated lower bound.  Finally use
\(\widehat B_{hw}(a)\ge y_{hw}\) and \eqref{eq:tiltedidentity}.
\end{proof}

\begin{proposition}[Master inequality]
Let a feasible depth-\(n\) tree have root mean \(m\).  For every integer
\(k\le n/2\),
\begin{equation}
 \boxed{
 \rho(a)^k
 \ge e^{-am}
 -2\Lambda(a)^{k/2}\sqrt{\frac{2amk}{n}}.
 }
 \label{eq:master}
\end{equation}
\end{proposition}

\begin{proof}
Put \(q=\lfloor n/k\rfloor\ge n/(2k)\).  The disjoint \(k\)-blocks telescope:
\[
 \sum_{j=0}^{q-1}\E D_{H_{jk}}^{(k)}
 =\E y_{H_{qk}}-y_\varnothing
 \le1-e^{-am}\le am.
\]
Hence one block start satisfies
\[
 \E D_{H_{jk}}^{(k)}\le\frac{2amk}{n}.
\]
Average \eqref{eq:blocklocal} over that starting history, use
\(\E m_{H_{jk}}=m\), Jensen, \(\widehat B_h\le1\), and
\(\E\sqrt D\le\sqrt{\E D}\).
\end{proof}

\begin{theorem}[Fixed-tilt logarithmic obstruction]
\label{thm:fixedtilt}
Assume \eqref{eq:equalmean}.  If \(\rho(a)<1\) for some fixed \(a>0\), set
\[
 d(a):=-\log\rho(a)>0,
 \qquad
 \Gamma(a):=2a+\frac{a\log\Lambda(a)}{d(a)}.
\]
Then
\begin{equation}
 \boxed{
 \liminf_{n\to\infty}\frac{\cB_n}{\log n}
 \ge\frac1{\Gamma(a)}.
 }
 \label{eq:fixedbound}
\end{equation}
Here and below \(\log\) denotes the natural logarithm.
\end{theorem}

\begin{proof}
For a feasible depth-\(n\) tree of root mean \(m\), choose
\[
 k=\left\lceil\frac{am+\log2}{d(a)}\right\rceil.
\]
Then \(\rho(a)^k\le\frac12e^{-am}\).  If \(k>n/2\), the definition of
\(k\) already gives a lower bound on \(m\) linear in \(n\), stronger than
\eqref{eq:fixedbound}.  Otherwise combine this estimate with
\eqref{eq:master} and square to obtain
\[
 n\le32amk\,\Lambda(a)^k e^{2am}.
\]
The ceiling bound on \(k\) therefore yields a finite constant \(C_a\) such that
\[
 n\le C_a(m+1)^2e^{\Gamma(a)m}.
\]
Taking logarithms,
\[
 \log n\le\Gamma(a)m+2\log(m+1)+\log C_a.
\]
Apply this inequality to a sequence of feasible trees whose root means are
within \(1/n\) of \(\cB_n\), and pass to a subsequence realizing the liminf.
If the ratio is finite, then \(m=O(\log n)\) on that subsequence and
\(\log(m+1)=o(\log n)\).  Division by \(\log n\) proves
\eqref{eq:fixedbound}.
\end{proof}

\begin{theorem}[Third-cumulant logarithmic obstruction]
\label{thm:third}
Assume that \(\mu_P\) and \(\mu_Q\) have common mean \(\mu\), common variance
\(\sigma^2>0\), and
\[
 \Delta\kappa_3:=\kappa_3(\mu_Q)-\kappa_3(\mu_P)>0.
\]
Then
\begin{equation}
 \boxed{
 \liminf_{n\to\infty}\frac{\cB_n}{\log n}
 \ge\frac{\Delta\kappa_3}{6\sigma^2}.
 }
 \label{eq:third}
\end{equation}
In particular \(\cB_n\to\infty\).
\end{theorem}

\begin{proof}
Let
\[
 K_P(a):=\log\E e^{-aX},
 \qquad
 K_Q(a):=\log\E e^{-aY}.
\]
Finite support gives, as \(a\downarrow0\),
\[
 K_P(a)=-\mu a+\frac{\sigma^2}{2}a^2-\frac{\kappa_3(P)}6a^3+O(a^4),
\]
\[
 K_Q(a)=-\mu a+\frac{\sigma^2}{2}a^2-\frac{\kappa_3(Q)}6a^3+O(a^4).
\]
Hence
\[
 d(a)=K_P(a)-K_Q(a)=\frac{\Delta\kappa_3}{6}a^3+O(a^4),
\]
and
\[
 \log\Lambda(a)=K_P(2a)-2K_P(a)=\sigma^2a^2+O(a^3).
\]
Therefore
\[
 \Gamma(a)\longrightarrow\frac{6\sigma^2}{\Delta\kappa_3}.
\]
Apply Theorem~\ref{thm:fixedtilt} for each sufficiently small fixed \(a>0\),
first take \(n\to\infty\), and only then send \(a\downarrow0\).
\end{proof}

\section{The dyadic parity pair}
\label{sec:dyadic}

We now specialize to
\begin{equation}
 \boxed{
 \mu_P=\frac14\delta_2+\frac34\delta_4,
 \qquad
 \mu_Q=\frac34\delta_3+\frac14\delta_5.
 }
 \label{eq:pair}
\end{equation}
Both laws have mean \(7/2\), variance \(3/4\), and centered third cumulants
\[
 \kappa_3(P)=-\frac34,
 \qquad
 \kappa_3(Q)=+\frac34.
\]
Thus \(\Delta\kappa_3=3/2\).

\begin{corollary}[The \(1/3\) lower law]
\label{cor:one-third}
For \eqref{eq:pair},
\begin{equation}
 \boxed{
 \liminf_{n\to\infty}\frac{\cB_n}{\log n}\ge\frac13.
 }
 \label{eq:onethird}
\end{equation}
Equivalently,
\[
 \cB_n\ge\left(\frac13-o(1)\right)\log n.
\]
For base-two logarithms the coefficient is \((\log2)/3\).
\end{corollary}

With \(z=e^{-a}\), the increment PGFs are
\[
 G_P(z)=\frac14z^2+\frac34z^4,
 \qquad
 G_Q(z)=\frac34z^3+\frac14z^5,
\]
and
\begin{equation}
 \boxed{
 G_P(z)-G_Q(z)=\frac14z^2(1-z)^3.
 }
 \label{eq:cubicPGF}
\end{equation}

\begin{corollary}[No finite-mean infinite zero-repair tree]
There is no infinite history-indexed zero-repair tree for \eqref{eq:pair} with
finite root mean.
\end{corollary}

\begin{proof}
Such a tree would restrict to a feasible depth-\(n\) tree for every \(n\),
forcing \(\cB_n\) to be uniformly bounded by its root mean, contrary to
\eqref{eq:onethird}.
\end{proof}

\section{Viability kernels and Bellman--Fenchel duality}
\label{sec:dual}

For the dyadic pair, write the probability generating function of a buffer law
as the same symbol.  Multiplying \eqref{eq:ZR} by four gives
\begin{equation}
 \boxed{
 B_{h0}(z)+3z^2B_{h1}(z)=A(z)B_h(z),
 \qquad A(z):=z(3+z^2).
 }
 \label{eq:PGFrec}
\end{equation}
Define
\[
 \cK_0:=\cP(\Z_{\ge0})
\]
and recursively
\begin{equation}
 \cK_{n+1}:=
 \left\{B\in\cP:
 \exists B_0,B_1\in\cK_n\text{ satisfying \eqref{eq:PGFrec}}\right\}.
 \label{eq:Kndef}
\end{equation}
Then
\[
 \cB_n=\inf_{B\in\cK_n}\E_BJ.
\]

\begin{proposition}[Convexity]
Every \(\cK_n\) is convex.
\end{proposition}

\begin{proof}
The base \(\cK_0\) is convex.  If \(B,B'\in\cK_{n+1}\) have feasible child
pairs \((B_0,B_1)\) and \((B'_0,B'_1)\), then linearity of
\eqref{eq:PGFrec} shows that every convex combination of \(B,B'\) has the
corresponding convex combinations of the children, which lie in \(\cK_n\).
\end{proof}

For a convex set \(K\) of probability laws and a test function \(j\), define
its lower support functional by
\begin{equation}
 \ell_K(j):=\inf_{B\in K}\langle j,B\rangle.
 \label{eq:lowerSupport}
\end{equation}
Write \(\ell_n:=\ell_{\cK_n}\).  Then
\[
 \cB_n=\ell_n(\id),
 \qquad \id(t)=t.
\]
On a finite state space, lower support is a complete coordinate for a closed
convex viability set.

\begin{proposition}[Lossless finite-dimensional support representation]
Let \(X\) be finite, \(\Delta_X\) the probability simplex, and
\(K\subseteq\Delta_X\) nonempty, closed and convex.  Then
\begin{equation}
 K=
 \bigcap_{j\in\R^X}
 \{b\in\Delta_X:\langle j,b\rangle\ge\ell_K(j)\}.
 \label{eq:recoverK}
\end{equation}
\end{proposition}

\begin{proof}
The forward inclusion is immediate.  If \(b\notin K\), finite-dimensional
separation gives \(j\) such that
\(\langle j,b\rangle<\inf_{x\in K}\langle j,x\rangle\).
\end{proof}

We now state the abstract dual recursion.  Let \(X,Y\) be finite sets,
\(K\subseteq\Delta_X\) nonempty compact convex, and
\(C,S_1,\ldots,S_r:\R^X\to\R^Y\) linear.  For positive weights \(p_i\)
with \(\sum_i p_i=1\), define
\[
 \Phi(K):=
 \left\{b\in\Delta_X:\exists b_i\in K,\quad
 Cb=\sum_i p_iS_ib_i\right\}.
\]

\begin{theorem}[Bellman--Fenchel duality]
\label{thm:BF}
If \(\Phi(K)\neq\varnothing\), then for every \(j\in\R^X\),
\begin{equation}
 \boxed{
 \ell_{\Phi(K)}(j)
 =
 \sup_{\lambda\in\R^Y}
 \left\{
 \min_{x\in X}(j_x-(C^*\lambda)_x)
 +\sum_i p_i\ell_K(S_i^*\lambda)
 \right\}.
 }
 \label{eq:BF}
\end{equation}
\end{theorem}

\begin{proof}
Consider the convex program minimizing \(\langle j,b\rangle\) over
\((b,b_1,\ldots,b_r)\in\Delta_X\times K^r\) subject to the displayed linear
equality.  With multiplier \(\lambda\), the Lagrangian is
\[
 \langle j-C^*\lambda,b\rangle
 +\sum_i p_i\langle S_i^*\lambda,b_i\rangle.
\]
Minimizing over the product set gives the right-hand side of \eqref{eq:BF}.
The primal feasible set is compact and convex and the constraints are linear,
so finite-dimensional strong duality follows from separation
\cite{Rockafellar1970}.
\end{proof}

For the dyadic model,
\[
 (C_Q^*\lambda)(t)=\frac34\lambda(t+3)+\frac14\lambda(t+5),
\]
\[
 (S_2^*\lambda)(t)=\lambda(t+2),
 \qquad
 (S_4^*\lambda)(t)=\lambda(t+4).
\]
Thus the formal countable-state Bellman transform is
\begin{equation}
\boxed{
\begin{aligned}
(\cD\ell)(j)=\sup_{\lambda}\Bigg\{&
\inf_{t\ge0}\left[j(t)-\frac34\lambda(t+3)-\frac14\lambda(t+5)\right]\\
&+\frac14\ell(\lambda(\cdot+2))
+\frac34\ell(\lambda(\cdot+4))
\Bigg\}.
\end{aligned}}
\label{eq:Ddyadic}
\end{equation}
On a finite exact truncation, \(\ell_{n+1}=\cD\ell_n\).  On the full
countable state space, the same expression is always a valid lower certificate
for each admissible \(\lambda\), but equality requires a separate no-gap
theorem in a suitable weighted function space.

\begin{proposition}[Weak countable-state recursion]
For every admissible \(\lambda\) for which the pairings are finite,
\begin{align}
\ell_{n+1}(j)\ge{}&
\inf_{t\ge0}\left[j(t)-\frac34\lambda(t+3)-\frac14\lambda(t+5)\right]
\notag\\
&+\frac14\ell_n(\lambda(\cdot+2))
+\frac34\ell_n(\lambda(\cdot+4)).
\label{eq:weakD}
\end{align}
\end{proposition}

\begin{proof}
For any feasible parent and children, add the zero pairing of \(\lambda\) with
the constraint.  Bound the parent contribution by its pointwise infimum and
each child contribution by the corresponding lower support, then take the
infimum over feasible triples.
\end{proof}

The operator is order preserving and additively homogeneous on every domain
where the expressions are finite:
\[
 \ell_1\le\ell_2\Rightarrow\cD\ell_1\le\cD\ell_2,
 \qquad
 \cD(\ell+c)=\cD\ell+c.
\]
Consequently, on any invariant class on which the uniform norm is finite,
\[
 \|\cD\ell_1-\cD\ell_2\|_\infty
 \le\|\ell_1-\ell_2\|_\infty.
\]
There is also the exact multiplier gauge \(\lambda\sim\lambda+c\), a direct
consequence of probability normalization.  These are the structural features
of a topical or Shapley-type operator; they motivate, but do not by themselves
justify, nonlinear Perron--Frobenius methods
\cite{LemmensNussbaum2012,Lins2023,AkianGaubertHochart2020}.

\section{Ordered histories and the phase quotient}
\label{sec:phase}

The scalar information spectrum forgets order.  The causal tree does not.  We
now make the difference algebraic.

Let
\[
 V_m:=\R^{\bits^m}
\]
with canonical basis \(e_w\), \(w\in\bits^m\).  For a binary word \(w\), let
\(\wt(w)\) denote its Hamming weight.  Define the scalarization map
\begin{equation}
 \boxed{
 \Pi_m:V_m\to\R^{m+1},
 \qquad
 \Pi_m e_w=e_{\wt(w)}.
 }
 \label{eq:Pim}
\end{equation}
For a signed measure on ordered histories, \(\Pi_m\) simply sums its masses
within equal-Hamming-weight layers.  Since the total information of either
binary source is an affine function of Hamming weight, this is exactly the
projection from ordered histories to scalar total-information profiles.

\begin{definition}[Ambient phase space]
\begin{equation}
 \boxed{W_m:=\ker\Pi_m.}
 \label{eq:Wm}
\end{equation}
Elements of \(W_m\) are signed rearrangements invisible to scalar total
information.
\end{definition}

\begin{theorem}[Dimension of hidden ordering information]
\label{thm:phasedim}
For every \(m\ge1\),
\begin{equation}
 \boxed{\dim W_m=2^m-(m+1).}
 \label{eq:phasedim}
\end{equation}
\end{theorem}

\begin{proof}
The map \(\Pi_m\) is surjective: for every \(0\le k\le m\), any word of
weight \(k\) maps to the \(k\)-th basis vector.  Hence
\(\operatorname{rank}\Pi_m=m+1\), and rank-nullity gives
\eqref{eq:phasedim}.
\end{proof}

Define normalized layer vectors
\[
 u_k^{(m)}:=\binom{m}{k}^{-1}\sum_{\wt(w)=k}e_w,
 \qquad 0\le k\le m,
\]
and let
\[
 U_m:=\Span\{u_0^{(m)},\ldots,u_m^{(m)}\}.
\]
Then \(\Pi_m|_{U_m}\) is an isomorphism and
\begin{equation}
 V_m=U_m\oplus W_m.
 \label{eq:UW}
\end{equation}
Thus \(U_m\) is a canonical visible sector and \(W_m\) the invisible phase
sector.

For \(m=2\),
\[
 u_0=e_{00},
 \qquad
 u_1=\frac12(e_{01}+e_{10}),
 \qquad
 u_2=e_{11},
\]
and
\begin{equation}
 W_2=\Span\left\{\theta\right\},
 \qquad
 \theta:=\frac12(e_{01}-e_{10}).
 \label{eq:theta2}
\end{equation}
This is the first nontrivial ordering phase: the words \(01\) and \(10\) have
the same scalar information but different causal pasts.

\subsection{Two-block phase coordinates for buffer PGFs}

Iterating \eqref{eq:PGFrec} twice gives
\begin{equation}
 \boxed{
 B_{00}+qB_{01}+qB_{10}+q^2B_{11}=A^2B,
 \qquad q:=3z^2.
 }
 \label{eq:2block}
\end{equation}
Put
\begin{equation}
 S:=\frac{B_{01}+B_{10}}2,
 \qquad
 \Theta:=\frac{B_{01}-B_{10}}2.
 \label{eq:Stheta}
\end{equation}
Then the coarse equation loses \(\Theta\):
\begin{equation}
 \boxed{B_{00}+2qS+q^2B_{11}=A^2B.}
 \label{eq:coarse2}
\end{equation}

\begin{proposition}[Exact two-block lifting criterion]
\label{prop:lift}
Let \(B,B_{00},S,B_{11}\) be probability PGFs satisfying
\eqref{eq:coarse2}.  They admit a genuine two-step zero-repair lift if and
only if there exists a signed series \(\Theta\) with \(\Theta(1)=0\) such that
\[
 B_{01}=S+\Theta,
 \qquad
 B_{10}=S-\Theta
\]
are probability PGFs and
\begin{equation}
 \boxed{
 B_0=\frac{B_{00}+q(S+\Theta)}{A},
 \qquad
 B_1=\frac{S-\Theta+qB_{11}}{A}
 }
 \label{eq:phaseLift}
\end{equation}
are probability PGFs.
\end{proposition}

\begin{proof}
Necessity follows by taking \(S,\Theta\) from an existing two-step tree and
solving the first-level recursions for \(B_0,B_1\).  Conversely, if the four
PGFs in the statement are nonnegative and normalized, the formulas imply
\[
 B_{00}+qB_{01}=AB_0,
 \qquad
 B_{10}+qB_{11}=AB_1,
\]
and adding \(B_0+qB_1=AB\) is equivalent to \eqref{eq:coarse2}.  Hence the
full two-step tree closes exactly.
\end{proof}

\begin{remark}
Proposition~\ref{prop:lift} is a change of coordinates, not a solution of the
positivity problem.  Its value is that it isolates the missing datum: the
coarse scalar equation determines the visible part, while causal feasibility
asks whether the corresponding phase fiber contains a point compatible with
all PGF inequalities.
\end{remark}

\section{Exact two-symbol causal phase geometry}
\label{sec:5over16}

The phase quotient is not merely formal.  At the first nontrivial block length
one can characterize \emph{exactly} which scalar joint profiles admit an
ordered causal lift.  This gives a closed-form phase-feasibility criterion,
not merely an LP description.

Let \(X=(X_1,X_2)\sim P^{\otimes2}\) and
\(W=(W_1,W_2)\sim Q^{\otimes2}\), where
\[
 P(0)=\frac14,\quad P(1)=\frac34,
 \qquad
 Q(0)=\frac34,\quad Q(1)=\frac14.
\]
For a coupling \(\pi\), define its scalar Hamming-layer profile
\begin{equation}
 r_{ij}:=\Pp_\pi\{\wt(X)=i,\ \wt(W)=j\},
 \qquad 0\le i,j\le2.
 \label{eq:rprofile}
\end{equation}
Its row and column marginals are fixed:
\[
 p^{(2)}=\frac1{16}(1,6,9),
 \qquad
 q^{(2)}=\frac1{16}(9,6,1).
\]
Write
\begin{equation}
 a:=r_{00},\qquad b:=r_{01},\qquad c:=r_{10},\qquad d:=r_{11},
 \qquad S:=a+b+c+d.
 \label{eq:abcd}
\end{equation}
Then every static scalar coupling is uniquely of the form
\begin{equation}
 \boxed{
 r(a,b,c,d)=
 \begin{pmatrix}
 a & b & \frac1{16}-a-b\\[1mm]
 c & d & \frac38-c-d\\[1mm]
 \frac9{16}-a-c & \frac38-b-d & S-\frac38
 \end{pmatrix}.}
 \label{eq:rstaticmatrix}
\end{equation}
Thus static feasibility is exactly the nonnegativity of the nine entries in
\eqref{eq:rstaticmatrix}.

A coupling is causal from \(W\) to \(X\) if the conditional law of \(X_1\)
given \((W_1,W_2)\) depends only on \(W_1\).  Equivalently,
\begin{equation}
 \sum_{x_2}\pi((a,x_2),(b,c))
 =Q(c)\sum_{x_2,w_2}\pi((a,x_2),(b,w_2))
 \label{eq:causalLPconstraints}
\end{equation}
for all \(a,b,c\in\bits\).

\begin{theorem}[Exact two-symbol causal scalar-profile criterion]
\label{thm:twofacet}
Let \(r=r(a,b,c,d)\) be a static scalar coupling with the fixed marginals
above.  Then \(r\) admits a causal lift \(\pi\) if and only if
\begin{align}
 \boxed{2a+3b+d\ge\frac3{16}},
 \label{eq:facet1}\\
 \boxed{2a+3b+2c+3d\le\frac98}.
 \label{eq:facet2}
\end{align}
Equivalently, the projected causal polytope is obtained from the static
transport polytope by intersecting it with precisely these two additional
half-spaces.

More constructively, a causal lift exists if and only if the interval
\begin{equation}
 \boxed{L(a,b,c,d)\le t\le U(a,b,c,d)}
 \label{eq:tinterval}
\end{equation}
is nonempty, where
\begin{align}
 L&:=\max\left\{
 0,\frac{4a}{3},4\left(S-\frac38\right),
 \frac38-2(b+d)\right\},
 \label{eq:Ldef}\\
 U&:=\min\left\{
 \frac14,4(a+b),\frac38-2b,\frac{4(a+c)}3\right\}.
 \label{eq:Udef}
\end{align}
\end{theorem}

\begin{proof}
Let
\[
 \alpha_{uv}:=\Pp_\pi(X_1=u,W_1=v).
\]
The one-symbol marginals force every such first-step coupling to have the form
\begin{equation}
 \alpha(t)=
 \begin{pmatrix}
 t&\frac14-t\\[1mm]
 \frac34-t&t
 \end{pmatrix},
 \qquad 0\le t\le\frac14.
 \label{eq:alphat}
\end{equation}
By causality and the product law of \(W\),
\begin{equation}
 \Pp_\pi(X_1=u,W_1=v,W_2=w)=Q(w)\alpha_{uv}(t).
 \label{eq:causalgroupmass}
\end{equation}
For \(u,v,w\in\bits\), put
\[
 s_{uvw}:=\pi((u,0),(v,w)).
\]
The companion mass with \(X_2=1\) is then
\begin{equation}
 \pi((u,1),(v,w))=Q(w)\alpha_{uv}(t)-s_{uvw}.
 \label{eq:companionmass}
\end{equation}
Hence a causal coupling is equivalent to choosing the eight numbers
\(s_{uvw}\) inside their elementary capacity intervals.

The scalar profile determines the following quantities:
\begin{align}
 s_{000}&=a,
 &s_{001}+s_{010}&=b,
 &s_{011}&=\frac1{16}-a-b,
 \label{eq:srow0}\\
 s_{100}&=a+c-\frac34t,
 &s_{101}+s_{110}&=b+d-\frac3{16}+\frac12t,
 &s_{111}&=\frac38-S+\frac14t.
 \label{eq:srow1}
\end{align}
Their capacity bounds from \eqref{eq:causalgroupmass} are
\begin{align*}
 0&\le s_{000}\le\frac34t,
 &0&\le s_{001}\le\frac14t,
 &0&\le s_{010}\le\frac3{16}-\frac34t,
 &0&\le s_{011}\le\frac1{16}-\frac14t,\\
 0&\le s_{100}\le\frac9{16}-\frac34t,
 &0&\le s_{101}\le\frac3{16}-\frac14t,
 &0&\le s_{110}\le\frac34t,
 &0&\le s_{111}\le\frac14t.
\end{align*}
For two nonnegative bins of capacities \(u,v\), a prescribed total \(R\)
can be split between them if and only if \(0\le R\le u+v\).  Applying this
observation to the two sums in \eqref{eq:srow0}--\eqref{eq:srow1}, and using
static feasibility for the bounds independent of \(t\), gives exactly
\eqref{eq:tinterval}--\eqref{eq:Udef}.

It remains to eliminate \(t\).  Write the four lower bounds in \(L\) as
\(L_0,L_1,L_2,L_3\) in the displayed order, and similarly
\(U_0,U_1,U_2,U_3\).  Nonemptiness is equivalent to
\(L_i\le U_j\) for every \(i,j\).  Under the static inequalities from
\eqref{eq:rstaticmatrix}, the only new comparisons are
\begin{equation}
 L_3\le U_1
 \quad\Longleftrightarrow\quad
 2a+3b+d\ge\frac3{16},
 \label{eq:compare1}
\end{equation}
and
\begin{equation}
 L_2\le U_3
 \quad\Longleftrightarrow\quad
 2a+3b+2c+3d\le\frac98.
 \label{eq:compare2}
\end{equation}
For completeness, the remaining potentially nontrivial comparisons follow as
follows.  From \(a+b\le1/16\) and \eqref{eq:facet1},
\(b+d\ge1/16\), hence \(L_3\le U_0\).  Also
\(S\le1/16+3/8=7/16\), which gives \(L_2\le U_0\); the inequality
\(L_2\le U_1\) is just \(c+d\le3/8\); and
\(L_2\le U_2\) follows from \(S\le7/16\) and \(b\le1/16\).
Finally, \(L_3\le U_3\) follows from \(S\ge3/8\), while comparisons involving
\(L_0,L_1\), together with \(L_3\le U_2\), are immediate from the static
nonnegativity constraints.  This proves that \eqref{eq:facet1}--\eqref{eq:facet2}
are necessary and sufficient.

For constructive sufficiency choose any \(t\in[L,U]\).  Set the four
singleton quantities in \eqref{eq:srow0}--\eqref{eq:srow1}.  Split \(b\)
between \(s_{001},s_{010}\) within their two capacities, and split
\(b+d-3/16+t/2\) between \(s_{101},s_{110}\) within theirs.  The interval
conditions guarantee that both splits exist.  Define every remaining cell by
\eqref{eq:companionmass}.  All sixteen masses are nonnegative and obey
\eqref{eq:causalLPconstraints}.  Equation \eqref{eq:causalgroupmass} gives
\(W\sim Q^{\otimes2}\).  The matrix \(\alpha(t)\) fixes the first-symbol
marginal of \(X\) to be \(P\); together with the Hamming-weight row
marginals \(p^{(2)}\), this forces the ordered word probabilities
\((1,3,3,9)/16\), hence \(X\sim P^{\otimes2}\).  By construction the
coupling projects to \(r\).  Thus it is the desired causal lift.
\end{proof}

Both facets are nonredundant.  The static profile with
\((a,b,c,d)=(0,0,3/8,0)\) violates \eqref{eq:facet1} while satisfying
\eqref{eq:facet2}; the static profile with
\((a,b,c,d)=(1/16,0,0,3/8)\) satisfies \eqref{eq:facet1} but violates
\eqref{eq:facet2}.

\begin{remark}[A first intrinsic phase criterion]
Theorem~\ref{thm:twofacet} is stronger than saying that a finite LP can be
solved.  At the first nontrivial block it eliminates the hidden ordering
variables completely: a static scalar profile has an ordered causal lift if
and only if two explicit oblique facets are satisfied.  The phase variable
\(t\) is recoverable from the interval \eqref{eq:tinterval}.  This is the
finite-block analogue of the kind of intrinsic positivity criterion one would
want to propagate under renormalization.
\end{remark}

The two-symbol total-information levels are
\[
 I_P(x)=4+2\wt(x),
 \qquad
 I_Q(w)=6+2\wt(w).
\]
Thus equal information occurs exactly for layer pairs \((i,j)=(1,0)\) and
\((2,1)\).  The matched mass of a scalar profile \eqref{eq:rstaticmatrix} is
\begin{equation}
 M_{\rm match}(r)=r_{10}+r_{21}
 =c+\frac38-b-d,
 \label{eq:matchedmass}
\end{equation}
so its mismatch is
\begin{equation}
 M_{\rm mis}(r)=\frac58+b+d-c.
 \label{eq:mismatchscalar}
\end{equation}

\begin{corollary}[Exact two-symbol causal mismatch]
\label{thm:5over16}
Among all causal couplings \(\pi\) with the prescribed product marginals,
\begin{equation}
 \boxed{
 \min_\pi \Pp_\pi\{I_P(X)\ne I_Q(W)\}=\frac5{16}.
 }
 \label{eq:5over16}
\end{equation}
Equivalently, the maximum causally matchable equal-information mass is
\(11/16\).  Without causality the minimum mismatch is \(1/4\), so the causal
surcharge is exactly \(1/16\).
\end{corollary}

\begin{proof}
For every causally liftable scalar profile, \eqref{eq:facet1} and
\(a+b\le1/16\) imply
\[
 \frac3{16}
 \le2a+3b+d
 \le\frac18+b+d,
\]
hence
\begin{equation}
 b+d\ge\frac1{16}.
 \label{eq:bdlower}
\end{equation}
Static feasibility also gives \(c+d\le3/8\).  Since \(d\ge0\),
\[
 c-b-d=(c+d)-(b+2d)
 \le\frac38-(b+d)
 \le\frac5{16}.
\]
Substitution into \eqref{eq:mismatchscalar} yields
\[
 M_{\rm mis}(r)\ge\frac58-\frac5{16}=\frac5{16}.
\]
Equality is attained at
\begin{equation}
 (a,b,c,d)=\left(0,\frac1{16},\frac38,0\right),
 \qquad
 r=\frac1{16}
 \begin{pmatrix}
 0&1&0\\
 6&0&0\\
 3&5&1
 \end{pmatrix}.
 \label{eq:optimalr}
\end{equation}
This profile satisfies both \eqref{eq:facet1} and \eqref{eq:facet2}, so
Theorem~\ref{thm:twofacet} supplies a causal lift.  Therefore the lower bound
is sharp.
\end{proof}

\begin{remark}[Which facet prices the surcharge]
The second causal facet \eqref{eq:facet2} is genuine for the full projected
profile geometry, but it is inactive at the optimizer of the mismatch
functional.  The exact \(1/16\) surcharge is already forced by the first
phase facet \eqref{eq:facet1}.  In this precise sense the local causal price is
a shadow of one hidden-ordering feasibility constraint.
\end{remark}

\subsection{Independent rational primal--dual certificate}

The closed-form proof above does not require a solver.  As an independent
check, the original sixteen-variable causal LP admits exact rational primal and
dual certificates.

Order rows and columns as \(00,01,10,11\).  The coupling
\begin{equation}
 \boxed{
 \pi=\frac1{16}
 \begin{pmatrix}
 0&1&0&0\\
 3&0&0&0\\
 3&0&0&0\\
 3&2&3&1
 \end{pmatrix}}
 \label{eq:primalmatrix}
\end{equation}
has the required marginals, satisfies \eqref{eq:causalLPconstraints}, projects
to \eqref{eq:optimalr}, and has mismatch \(5/16\).

Vectorize \(\pi\) row by row and write all four \(X\)-marginals, four
\(W\)-marginals, and eight causal equations as \(A\pi=b\), where
\begin{equation}
 b=\frac1{16}(1,3,3,9,9,3,3,1,0,0,0,0,0,0,0,0)^T.
 \label{eq:5over16b}
\end{equation}
The dual vector
\begin{equation}
 y=
 \left(
 1,0,-\frac14,\frac34;
 0,0,0,-2;
 0,0,0,0,1,0,-3,0
 \right)
 \label{eq:dualvector}
\end{equation}
satisfies \(b^Ty=5/16\) and
\[
 A^Ty=
 \begin{pmatrix}
 1&1&1&-1\\
 0&0&0&-2\\
 0&-1&-1&0\\
 1&0&0&1
 \end{pmatrix}
 \le
 C:=
 \begin{pmatrix}
 1&1&1&1\\
 0&1&1&1\\
 0&1&1&1\\
 1&0&0&1
 \end{pmatrix}
\]
entrywise.  Hence weak duality independently certifies the same optimum.  The
complete rational matrix \(A\) and the multiplication are displayed in
Appendix~\ref{app:5over16matrix}.

The scalar two-symbol laws are
\begin{equation}
 \mu_P^{*2}=\frac1{16}\delta_4+\frac6{16}\delta_6+\frac9{16}\delta_8,
 \label{eq:P2}
\end{equation}
\begin{equation}
 \mu_Q^{*2}=\frac9{16}\delta_6+\frac6{16}\delta_8+\frac1{16}\delta_{10}.
 \label{eq:Q2}
\end{equation}
Static matching pairs \(6/16\) at level \(6\) and \(6/16\) at level \(8\),
for common mass \(12/16\).  Hence
\[
 R_P^{\rm stat}=\frac1{16}\delta_4+\frac3{16}\delta_8,
 \qquad
 R_Q^{\rm stat}=\frac3{16}\delta_6+\frac1{16}\delta_{10}.
\]
The causal optimizer \eqref{eq:primalmatrix} matches \(6/16\) at level \(6\)
but only \(5/16\) at level \(8\), so
\[
 R_P^{\rm causal}=\frac1{16}\delta_4+\frac4{16}\delta_8,
\]
\[
 R_Q^{\rm causal}=\frac3{16}\delta_6+\frac1{16}\delta_8
 +\frac1{16}\delta_{10}.
\]
Consequently,
\begin{equation}
 \boxed{
 R_P^{\rm causal}=R_P^{\rm stat}+\frac1{16}\delta_8,
 \qquad
 R_Q^{\rm causal}=R_Q^{\rm stat}+\frac1{16}\delta_8.}
 \label{eq:hiddenres}
\end{equation}
and therefore
\begin{equation}
 \boxed{
 R_P^{\rm causal}-R_Q^{\rm causal}
 =R_P^{\rm stat}-R_Q^{\rm stat}.}
 \label{eq:sameSigned}
\end{equation}

\begin{corollary}[A causal obstruction invisible to signed scalar spectra]
The additional causal mismatch \(1/16\) in Corollary~\ref{thm:5over16}
leaves the signed scalar residual profile unchanged.  Hence the signed
information-spectrum difference alone cannot detect all causal matching
obstructions.
\end{corollary}

The first nontrivial causal phase obstruction is therefore fully visible in
three equivalent forms: an ordered-history feasibility problem, the two
intrinsic scalar facets \eqref{eq:facet1}--\eqref{eq:facet2}, and an exact
primal--dual LP certificate.

\section{Phase creation under block doubling}
\label{sec:doublingphase}

The one-dimensional phase \(e_{01}-e_{10}\) is special to two symbols.  At
larger scales the hidden sector grows rapidly, and block doubling creates new
phase directions even if the input blocks are individually scalarized.

Identify
\[
 V_{2m}\cong V_m\otimes V_m
\]
by concatenating two length-\(m\) words.  Define
\[
 C_m:\R^{m+1}\otimes\R^{m+1}\to\R^{2m+1}
\]
on basis vectors by
\begin{equation}
 C_m(e_i\otimes e_j)=e_{i+j}.
 \label{eq:Cm}
\end{equation}
Then the scalarization maps satisfy the exact factorization
\begin{equation}
 \boxed{
 \Pi_{2m}=C_m\circ(\Pi_m\otimes\Pi_m).
 }
 \label{eq:Pifactor}
\end{equation}

Let
\begin{equation}
 K_m:=\ker\left(C_m|_{\Pi_m(U_m)\otimes\Pi_m(U_m)}\right),
 \label{eq:Km}
\end{equation}
and identify this kernel back inside \(U_m\otimes U_m\) using the isomorphism
\(\Pi_m|_{U_m}\).

\begin{theorem}[Exact phase decomposition under doubling]
\label{thm:doublingdecomp}
For every \(m\ge1\),
\begin{equation}
\boxed{
W_{2m}
=(W_m\otimes W_m)
\oplus(W_m\otimes U_m)
\oplus(U_m\otimes W_m)
\oplus K_m.
}
\label{eq:phaseDecomp}
\end{equation}
Moreover,
\begin{equation}
 \boxed{\dim K_m=m^2.}
 \label{eq:Kdim}
\end{equation}
\end{theorem}

\begin{proof}
From \eqref{eq:UW},
\[
V_m\otimes V_m
=(U_m\otimes U_m)\oplus(U_m\otimes W_m)
\oplus(W_m\otimes U_m)\oplus(W_m\otimes W_m).
\]
Equation \eqref{eq:Pifactor} annihilates every tensor component containing a
factor from \(W_m\).  On \(U_m\otimes U_m\), it is identified with \(C_m\).
Therefore the kernel is exactly the direct sum in \eqref{eq:phaseDecomp}.

The map \(C_m\) is surjective because every basis vector \(e_k\),
\(0\le k\le2m\), equals \(C_m(e_i\otimes e_j)\) for some
\(0\le i,j\le m\) with \(i+j=k\).  Hence
\[
 \dim K_m=(m+1)^2-(2m+1)=m^2.
\]
\end{proof}

The sector \(K_m\) has a simple interpretation.  Even if each half-block has
already been reduced to its Hamming layer, the total layer \(i+j\) forgets how
the weight was split between the two halves.  Thus \(K_m\) is a genuinely
\emph{fresh anti-diagonal phase sector}: it is created by coarse-graining itself,
not inherited from unresolved phase inside either sub-block.

\subsection{The first doubling: two symbols to four}

For \(m=2\), \(\dim W_2=1\) and \(\dim U_2=3\).  Hence
\[
 \dim W_4=1+3+3+4=11,
\]
in agreement with \(2^4-5=11\).  Seven dimensions are inherited:
\[
 \theta\otimes\theta,
 \qquad
 \theta\otimes u_i,
 \qquad
 u_i\otimes\theta,
 \qquad i=0,1,2.
\]
A convenient basis of the four-dimensional fresh sector \(K_2\) is
\begin{align*}
\kappa_1&=u_0\otimes u_1-u_1\otimes u_0,\\
\kappa_2&=u_0\otimes u_2-u_2\otimes u_0,\\
\kappa_3&=u_1\otimes u_2-u_2\otimes u_1,\\
\kappa_4&=u_0\otimes u_2+u_2\otimes u_0-2u_1\otimes u_1.
\end{align*}
Each vector has zero anti-diagonal sum under \(C_2\).

\begin{corollary}[No globally one-dimensional phase model]
The phase direction \(01-10\) cannot close under block doubling.  Any exact
scale description retaining all ordering information must allow additional
phase directions, including a fresh \(m^2\)-dimensional sector at scale
\(m\to2m\).
\end{corollary}

This corollary is kinematic: it concerns the ambient phase space.  The local
fixed-block subset selected by positivity and causality will be characterized
exactly in Section~\ref{sec:causal-farkas}.  The genuinely open object is
smaller still: the subset whose descendants remain zero-repair viable across
successive scales.

\section{The cubic defect as a spectral normal form}
\label{sec:cubic}

Let \(E\) be the forward shift on test functions,
\[
 (Ef)(t):=f(t+1).
\]
For the dyadic pair define
\[
 P^*:=\frac14E^2+\frac34E^4,
 \qquad
 Q^*:=\frac34E^3+\frac14E^5.
\]

\begin{theorem}[Exact cubic shift factorization]
\label{thm:cubicfactor}
\begin{equation}
 \boxed{
 P^*-Q^*=\frac14E^2(I-E)^3.
 }
 \label{eq:shiftfactor}
\end{equation}
Equivalently,
\begin{equation}
 \boxed{
 \mu_P-\mu_Q
 =\frac14\delta_2*(\delta_0-\delta_1)^{*3}.
 }
 \label{eq:measurefactor}
\end{equation}
\end{theorem}

\begin{proof}
Expand \(\frac14E^2(I-E)^3\).  The coefficients are
\(\frac14(1,-3,3,-1)\) on shifts \(2,3,4,5\), exactly those of
\(P^*-Q^*\).
\end{proof}

Thus every polynomial of degree at most two is annihilated.  On the cubic,
\begin{equation}
 (P^*-Q^*)t^3=-\frac32.
 \label{eq:cubicpoly}
\end{equation}
The first transverse polynomial jet is therefore one-dimensional modulo
quadratics.

More generally, if a finite signed measure \(\nu=\sum_k\nu_k\delta_k\) has
vanishing raw moments through order \(r-1\), its generating polynomial
\(V(z)=\sum_k\nu_kz^k\) is divisible by \((1-z)^r\).  This is simply the
invertible triangular relation between raw moments and falling-factorial
moments.  On test functions, the associated difference operator factors by
\((I-E)^r\).

Laplace characters
\[
 f_a(t):=e^{-at}
\]
are eigenfunctions of the shift, \(Ef_a=e^{-a}f_a\).  Therefore
\begin{equation}
 \boxed{
 (P^*-Q^*)f_a
 =\frac14e^{-2a}(1-e^{-a})^3f_a.
 }
 \label{eq:laplaceSymbol}
\end{equation}
The transform used in Section~\ref{sec:lower} is thus a spectral character of
the exact translation algebra, not an unrelated trick.

\begin{lemma}[Shift invariance of viability]
If \(B\in\cK_n\), then \(z^sB\in\cK_n\) for every integer \(s\ge0\).
\end{lemma}

\begin{proof}
Multiply \eqref{eq:PGFrec} at every node of a feasible tree by \(z^s\).
\end{proof}

\begin{remark}[Why Laplace modes are multipliers, not bill coordinates]
If \(a>0\) and \(\cK_n\neq\varnothing\), then
\[
 \ell_n(e^{-a\cdot})=0.
\]
Indeed, shift any fixed feasible law by \(s\): its pairing with
\(e^{-a\cdot}\) is multiplied by \(e^{-as}\) and therefore tends to zero as
\(s\to\infty\); nonnegativity gives the reverse inequality.  Laplace modes
are therefore useful as dual multipliers and shift eigencharacters, not as
direct coordinates of \(\cB_n=\ell_n(\id)\).
\end{remark}

\subsection{Why the third order is critical for this converse}

Suppose more generally that
\[
 d(a):=-\log\rho(a)=c_ra^r+O(a^{r+1}),
 \qquad c_r>0,
\]
while \(\log\Lambda(a)=\sigma^2a^2+O(a^3)\).  Then
\begin{equation}
 \Gamma(a)=2a+\frac{\sigma^2}{c_r}a^{3-r}(1+o(1)).
 \label{eq:criticality}
\end{equation}
For this proof mechanism, \(r=3\) is critical: \(r=2\) makes the logarithmic
scale too small, \(r=3\) gives a finite nonzero coefficient, and \(r>3\)
requires a different scaling or a sharper dual.  This is a statement about the
present converse mechanism, not a universal exponent theorem.

\section{Green coordinates and boundary renormalization}
\label{sec:green}

Set
\[
 P(z):=1+3z^2,
 \qquad
 A(z):=z(3+z^2).
\]
Then
\begin{equation}
 \boxed{P(z)-A(z)=(1-z)^3.}
 \label{eq:cubiczero}
\end{equation}
The original PGFs are regular at \(z=1\); the analytic difficulty is the
triple zero.  We resolve it by dividing by the corresponding Green kernel.

\begin{definition}[Cubic Green potential]
For a buffer PGF \(B\), define
\begin{equation}
 \boxed{
 T_B(z):=\frac{1-B(z)}{(1-z)^3}.
 }
 \label{eq:TB}
\end{equation}
\end{definition}

\begin{theorem}[Exact affine Green conjugacy]
The one-step zero-repair recursion \eqref{eq:PGFrec} is equivalent to
\begin{equation}
 \boxed{
 T_{B_0}(z)+3z^2T_{B_1}(z)=A(z)T_B(z)+1.
 }
 \label{eq:GreenRec}
\end{equation}
\end{theorem}

\begin{proof}
Using \(P-A=(1-z)^3\),
\begin{align*}
(1-B_0)+3z^2(1-B_1)
&=P-AB\\
&=A(1-B)+(P-A).
\end{align*}
Divide by \((1-z)^3\).
\end{proof}

Thus the signed cubic source becomes the positive unit source.  Positivity of
the probability cone, rather than inversion of the cubic difference, is the
remaining barrier.

Let \(\beta_n:=\Pp_B(J>n)\) and
\(T_B(z)=\sum_{n\ge0}t_nz^n\).  Direct coefficient extraction gives
\begin{equation}
 t_n=\sum_{k=0}^n(n-k+1)\beta_k.
 \label{eq:tcoeff}
\end{equation}
With the backward difference \(\nabla t_n=t_n-t_{n-1}\),
\begin{equation}
 \boxed{
 \nabla t_n=\sum_{k=0}^n\beta_k,
 \qquad
 \nabla^2t_n=\beta_n,
 \qquad
 \nabla^3t_n=-\Pp(J=n)\ (n\ge1).
 }
 \label{eq:curvature}
\end{equation}
In particular,
\begin{equation}
 \boxed{
 \E_BJ=\sum_{n\ge0}\beta_n
 =\lim_{n\to\infty}\nabla t_n.
 }
 \label{eq:meanSlope}
\end{equation}
The bill is therefore an asymptotic slope in Green coordinates.

Iterating \eqref{eq:GreenRec} for \(m\) generations gives
\begin{equation}
 \sum_{w\in\bits^m}3^{\wt(w)}z^{2\wt(w)}T_{B_{hw}}(z)
 =A(z)^mT_{B_h}(z)+G_m(z),
 \label{eq:blockGreen}
\end{equation}
where
\begin{equation}
 \boxed{
 G_m(z)=\frac{P(z)^m-A(z)^m}{(1-z)^3}
 =\sum_{j=0}^{m-1}P(z)^{m-1-j}A(z)^j.
 }
 \label{eq:Gm}
\end{equation}
Hence \(G_m\) is coefficientwise nonnegative.  At \(m=2\),
\begin{equation}
 \boxed{G_2(z)=P(z)+A(z)=(1+z)^3.}
 \label{eq:G2}
\end{equation}
The normalized mask \((1+z)^3/4\) is the standard binary quadratic B-spline
subdivision mask (B-spline order three).  The connection is exact at the
algebraic level: subdivision theory relates such masks, vanishing moments and
polynomial reproduction \cite{CharinaConti2013}.

Define the positive boundary dilation
\begin{equation}
 \boxed{
 (\cR F)(z):=\frac14(1+z)^3F(z^2).
 }
 \label{eq:Rboundary}
\end{equation}
This is not asserted to be the full nonlinear Bellman renormalization.  It is
the exact scale map in the cubic Green boundary algebra.

\begin{proposition}[Exact third-difference conjugacy]
Let \(\Delta_3F:=(1-z)^3F\) and \((D_2G)(z):=G(z^2)\).  Then
\begin{equation}
 \boxed{
 \Delta_3\cR=\frac14D_2\Delta_3.
 }
 \label{eq:diffconj}
\end{equation}
\end{proposition}

\begin{proof}
Use \((1-z)^3(1+z)^3=(1-z^2)^3\).
\end{proof}

\begin{proposition}[Boundary pole spectrum]
If
\[
 F(z)=\frac{c+o(1)}{(1-z)^r}
 \qquad(z\uparrow1),
\]
then
\begin{equation}
 \boxed{
 (\cR F)(z)=2^{1-r}\frac{c+o(1)}{(1-z)^r}.
 }
 \label{eq:polespectrum}
\end{equation}
\end{proposition}

Thus pole orders \(r=0,1,2\) have multipliers \(2,1,1/2\), respectively.
The word ``spectrum'' here refers only to the exact boundary scaling law
\eqref{eq:polespectrum}; no spectrum of the full viability operator has yet
been proved.

\section{An explicitly controllable marginal phase channel}
\label{sec:eta}

A two-block causal correction arising in the dyadic model is the signed packet
\begin{equation}
 \boxed{
 \eta(z)=\frac1{16}(z^2+z^4-1-z^6)
 =-\frac1{16}(1-z^2)(1-z^4).
 }
 \label{eq:eta}
\end{equation}
It is mass- and mean-neutral:
\[
 \eta(1)=0,
 \qquad
 \eta'(1)=0,
\]
while
\begin{equation}
 \eta''(1)=-1.
 \label{eq:eta2}
\end{equation}
Under the signed residual dilation
\[
 (\cS\Theta)(z):=\frac14\Theta(z^2),
\]
the factorial moments \(M_j=\Theta^{(j)}(1)\) transform as
\begin{align*}
M_0&\mapsto\frac14M_0,\\
M_1&\mapsto\frac12M_1,\\
M_2&\mapsto M_2+\frac12M_1,\\
M_3&\mapsto2M_3+3M_2.
\end{align*}
Hence \(M_2\) is exactly marginal on the subspace \(M_0=M_1=0\).

Apply the cubic Green transform to \(\eta\):
\begin{equation}
 \boxed{
 \Phi(z):=-\frac{\eta(z)}{(1-z)^3}
 =\frac{(1+z)^2(1+z^2)}{16(1-z)}.
 }
 \label{eq:Phi}
\end{equation}
Its coefficients are
\[
 [z^n]\Phi=\frac1{16}(1,3,5,7,8,8,8,\ldots),
\]
so \(\Phi\) lies in the simple-pole marginal sector of
\eqref{eq:polespectrum} and carries zero asymptotic mean slope.

\begin{proposition}[Positive cohomological defect]
\begin{equation}
 \boxed{
 \Phi-\cR\Phi
 =\frac{(1+z)^3(1+z^2)(z^4+2z^2+3)}{64}.
 }
 \label{eq:cohomology}
\end{equation}
The right-hand side has nonnegative coefficients.  Consequently,
\begin{equation}
 0\le\cR^k\Phi\le\Phi
 \qquad(k\ge0)
 \label{eq:Rbound}
\end{equation}
coefficientwise.
\end{proposition}

\begin{proof}
Direct substitution into \eqref{eq:Rboundary} gives
\eqref{eq:cohomology}.  Positivity of \(\cR\) then allows iteration.
\end{proof}

Define
\[
 \eta_s(z):=4^{-s}\eta(z^{2^s}),
 \qquad s\ge0,
\]
and
\begin{equation}
 \boxed{
 E_*(z):=\frac23z^2+
 \sum_{s\ge0}\frac{4^{-s}}{16}
 \left(1+z^{2^{s+1}}+z^{2^{s+2}}+z^{6\cdot2^s}\right).
 }
 \label{eq:envelope}
\end{equation}

\begin{theorem}[Universal finite-mean envelope for the explicit phase channel]
\label{thm:envelope}
The series \(E_*\) is a probability PGF with
\begin{equation}
 \boxed{\E_{E_*}J=\frac{17}{6}.}
 \label{eq:17over6}
\end{equation}
For every sequence \((\theta_s)_{s\ge0}\) with \(|\theta_s|\le1\),
\begin{equation}
 E_*(z)+\sum_{s\ge0}\theta_s\eta_s(z)
 \label{eq:perturbedE}
\end{equation}
is again a probability PGF with the same mean \(17/6\).
\end{theorem}

\begin{proof}
The envelope part has total mass
\[
 \sum_{s\ge0}\frac{4\cdot4^{-s}}{16}=\frac13,
\]
which together with \(2/3\) from the first atom gives mass one.  The first
term contributes mean \(4/3\), while scale \(s\) contributes
\[
 \frac{4^{-s}}{16}
 \left(2^{s+1}+2^{s+2}+6\cdot2^s\right)=\frac34\,2^{-s}.
\]
Summing yields \(4/3+3/2=17/6\).  Coefficientwise, \(E_*\) dominates the
sum of absolute masses of every packet that can contribute at that
coefficient; the triangle inequality gives nonnegativity of
\eqref{eq:perturbedE}.  Each \(\eta_s\) has zero mass and zero first moment,
so normalization and mean are unchanged.
\end{proof}

\begin{remark}[The limit of this theorem]
Theorem~\ref{thm:envelope} controls one explicit phase channel through
arbitrarily many scales at \(O(1)\) mean cost.  It does not show that the full
causal phase fiber decomposes into copies of this channel.  Theorem
\ref{thm:doublingdecomp} in fact warns that fresh phase dimensions are created
at every blocking step.
\end{remark}

\section{A critical self-similar reservoir}
\label{sec:mahler}

The dilation algebra also contains a canonical Mahler fixed point.  Define
\begin{equation}
 (\mathfrak M B)(z):=\frac34z^2+\frac14B(z^2).
 \label{eq:MahlerMap}
\end{equation}

\begin{proposition}[Mahler fixed point]
There is a unique probability PGF satisfying \(B=\mathfrak M B\), namely
\begin{equation}
 \boxed{
 B_*(z)=\sum_{k\ge1}\frac3{4^k}z^{2^k}.
 }
 \label{eq:MahlerFixed}
\end{equation}
It has
\begin{equation}
 \E_{B_*}J=3,
 \qquad
 \E_{B_*}J^2=\infty.
 \label{eq:MahlerMoments}
\end{equation}
\end{proposition}

\begin{proof}
Iterating \(B(z)=\frac34z^2+\frac14B(z^2)\) determines the coefficients and
gives \eqref{eq:MahlerFixed}; its masses sum to one.  The first moment is
\[
 \sum_{k\ge1}\frac3{4^k}2^k=3,
\]
while the second moment is
\[
 \sum_{k\ge1}\frac3{4^k}4^k=\infty.
\]
\end{proof}

The fixed point is therefore exactly critical between finite first and second
moments.  This fact should not be overinterpreted: a fixed point of the static
scale map is not automatically an infinite zero-repair tree, and
Theorem~\ref{thm:third} forbids any such tree with finite root mean.  The
Mahler law is best viewed as a canonical boundary object of the scale algebra,
not as a completed causal construction.

% ============================================================
\section{Positive fibers over scalar information}
\label{sec:positive-fiber}
% ============================================================

The quotient $V_m\to\R^{m+1}$ identifies exactly which linear information is
lost when one passes from ordered histories to total information.  Positivity
can also be described exactly before causality is imposed.

For $0\le k\le m$ recall the normalized Hamming-layer vector
\[
 u_k^{(m)}
 :=
 \frac{1}{\binom{m}{k}}
 \sum_{|w|=k}e_w,
\]
and define the canonical section
\begin{equation}
 \sigma_m(r)
 :=
 \sum_{k=0}^m r_k u_k^{(m)},
 \qquad
 r=(r_0,\ldots,r_m)\in\Delta_m.
 \label{eq:canonicalsection}
\end{equation}
Here $\Delta_m$ denotes the probability simplex on the $m+1$ Hamming
layers.

\begin{theorem}[Exact positive phase fiber]
\label{thm:positivefiber}
For $r\in\Delta_m$, the set of probability laws on ordered histories with
scalar profile $r$ is
\begin{equation}
 \boxed{
 \Pi_m^{-1}(r)\cap\Delta(\{0,1\}^m)
 =
 \sigma_m(r)+\mathsf F_m(r),
 }
 \label{eq:positivefiber}
\end{equation}
where
\begin{equation}
 \boxed{
 \mathsf F_m(r)
 =
 \left\{
 \theta\in W_m:
 \theta_w\ge
 -\frac{r_{|w|}}{\binom{m}{|w|}}
 \ \text{for every }w
 \right\}.
 }
 \label{eq:Fmr}
\end{equation}
Equivalently, the fiber is the direct product, over Hamming layers, of
scaled simplices of total mass $r_k$.  If every $r_k>0$, its affine dimension
is
\[
 \sum_{k=0}^m\left(\binom{m}{k}-1\right)
 =2^m-(m+1).
\]
\end{theorem}

\begin{proof}
If $\nu$ is a probability vector with $\Pi_m\nu=r$, then the total mass of
$\nu$ on the layer $|w|=k$ is $r_k$.  Since $\sigma_m(r)$ places the uniform
mass $r_k/\binom{m}{k}$ on that layer, the difference
$\theta=\nu-\sigma_m(r)$ has zero sum on each layer, hence
$\theta\in\ker\Pi_m=W_m$.  The condition $\nu_w\ge0$ is exactly
\eqref{eq:Fmr}.  The converse is immediate.  The product-of-simplices
statement follows layer by layer, and the dimension formula follows by
summing the dimensions $\binom{m}{k}-1$.
\end{proof}

Theorem~\ref{thm:positivefiber} is the exact analogue of a coefficientwise
positivity test at the \emph{ambient} history level.  It is not yet the causal
criterion we ultimately need.  Prefix compatibility cuts a generally much
smaller subset out of this product of simplices.  The two-symbol theorem in
Section~\ref{sec:5over16} proves that this cut is already strict at the
first nontrivial block length.

% ============================================================
\section{Exact finite-block causal phase fibers}
\label{sec:causal-farkas}
% ============================================================

At a fixed block length the causal restriction can be characterized without an
ansatz.  Let \(\Omega_m:=\bits^m\) and
\[
 \mathcal J_m:=\R^{\Omega_m\times\Omega_m}\cong V_m\otimes V_m.
\]
Define the joint scalarization
\begin{equation}
 \boxed{\Sigma_m:=\Pi_m\otimes\Pi_m,\qquad
 \Sigma_m e_{(x,w)}=e_{(\wt(x),\wt(w))}.}
 \label{eq:Sigmam}
\end{equation}
Since \(\Sigma_m\) is surjective,
\begin{equation}
 \boxed{\dim\ker\Sigma_m=4^m-(m+1)^2.}
 \label{eq:jointphase-dim}
\end{equation}
Thus even before causality is imposed, scalar joint profiles discard an
exponentially growing family of ordering directions.

For a scalar joint profile \(r=(r_{ij})_{0\le i,j\le m}\), put
\begin{equation}
 \boxed{\tau_m(r):=\sum_{i,j=0}^m r_{ij}\,u_i^{(m)}\otimes u_j^{(m)}.}
 \label{eq:taum}
\end{equation}
Then \(\Sigma_m\tau_m(r)=r\).

Let \(\pi(x,w)\) be a nonnegative coupling of \(P^{\otimes m}\) and
\(Q^{\otimes m}\).  Besides the two marginal families, causality from \(W\)
to \(X\) is equivalent to
\begin{equation}
 \boxed{
 \sum_{x'\in\Omega_{m-t}}\pi((a,x'),(b,c))
 =Q^{\otimes(m-t)}(c)
 \sum_{x'\in\Omega_{m-t}}\sum_{c'\in\Omega_{m-t}}
 \pi((a,x'),(b,c'))}
 \label{eq:generalcausality}
\end{equation}
for \(1\le t<m\), \(a,b\in\Omega_t\), and \(c\in\Omega_{m-t}\).  These
equations say exactly that, conditional on \(W_{1:t}=b\), the law of
\(X_{1:t}\) is independent of the future \(W_{t+1:m}\).  Write all marginal
and causal equations as
\begin{equation}
 A_m\pi=b_m.
 \label{eq:Am}
\end{equation}
Redundant rows are allowed.

\begin{definition}[Local causal phase fiber]
For a scalar joint profile \(r\), define
\begin{equation}
 \boxed{\mathsf C_m^{\mathrm{loc}}(r)
 :=\{\theta\in\ker\Sigma_m:
 \tau_m(r)+\theta\ge0,\;A_m(\tau_m(r)+\theta)=b_m\}.}
 \label{eq:Cloc}
\end{equation}
\end{definition}

\begin{theorem}[Exact finite-block causal phase-fiber criterion]
\label{thm:causal-farkas}
For every \(m\) and every scalar joint profile \(r\),
\begin{equation}
 \boxed{\mathsf C_m^{\mathrm{loc}}(r)=
 \{\theta\in\ker\Sigma_m:
 A_m\theta=b_m-A_m\tau_m(r),\;\theta\ge-\tau_m(r)\}.}
 \label{eq:ClocExact}
\end{equation}
Hence the local causal phase fiber is a bounded polytope, possibly empty.
If \(P,Q,r\) are rational, it is a rational polytope.

Moreover, the fiber is nonempty if and only if
\begin{equation}
 \boxed{\langle y,b_m\rangle+\langle z,r\rangle\ge0}
 \label{eq:FarkasCriterion}
\end{equation}
for every pair \((y,z)\) satisfying
\begin{equation}
 \boxed{A_m^Ty+\Sigma_m^Tz\ge0}
 \label{eq:FarkasCone}
\end{equation}
componentwise.  For rational \(P,Q\), the cone in
\eqref{eq:FarkasCone} is rational polyhedral, hence a finite rational
generating family of rays suffices.
\end{theorem}

\begin{proof}
Every coupling with scalar profile \(r\) has the unique form
\(\pi=\tau_m(r)+\theta\) with \(\theta\in\ker\Sigma_m\).  Substitution gives
\eqref{eq:ClocExact}; boundedness follows because the translated fiber lies in
the probability simplex.

For nonemptiness apply Farkas' lemma to
\[
 \pi\ge0,\qquad A_m\pi=b_m,\qquad\Sigma_m\pi=r.
\]
Such a nonnegative solution exists if and only if every \((y,z)\) with
\(A_m^Ty+\Sigma_m^Tz\ge0\) satisfies
\(\langle y,b_m\rangle+\langle z,r\rangle\ge0\).  Rational polyhedrality is
immediate when the source probabilities are rational.
\end{proof}

\begin{remark}[The real remaining difficulty]
Theorem~\ref{thm:causal-farkas} completely settles the \emph{local fixed-block}
causal fiber.  It does not settle zero repair.  A zero-repair block must land
in child laws that remain viable for all remaining depths.  Thus the hard
question is no longer whether a finite inequality description exists at fixed
\(m\), but whether those certificates admit a scale-stable intrinsic recursion
with uniform bill control under \(m\mapsto2m\).
\end{remark}

\begin{corollary}[Finite truncations admit exact certificates]
Every finite-horizon zero-repair problem with a finite support cutoff is a
finite linear feasibility problem.  Any projected coarse/phase fiber is
therefore a rational polytope and admits a finite exact Farkas certificate
family.  Numerical LP duals can consequently be rationalized into exact
certificates; what is not automatic is a support-independent, scale-uniform
formula.
\end{corollary}

% ============================================================
\section{Why finite closures cannot be the answer}
\label{sec:nogo}
% ============================================================

The phase-space growth above might tempt one to search for a fortunate finite
collection of exact buffer profiles or for one scalar potential that absorbs
all of the ordering information.  Both possibilities are ruled out.

\begin{theorem}[No finite exact profile automaton]
\label{thm:nofiniteautomaton}
There do not exist finitely many probability PGFs
$B^{(1)},\ldots,B^{(K)}$ and transition maps
$\tau_0,\tau_1:\{1,\ldots,K\}\to\{1,\ldots,K\}$ such that
\begin{equation}
 z(3+z^2)B^{(s)}(z)
 =B^{(\tau_0(s))}(z)+3z^2B^{(\tau_1(s))}(z)
 \label{eq:finiteautomaton}
\end{equation}
for every $s$.
\end{theorem}

\begin{proof}
Let $G=(B^{(1)},\ldots,B^{(K)})^\top$, and let $L,H$ be the selector
matrices induced by $\tau_0,\tau_1$.  Then
\[
 M(z)G(z)=0,
 \qquad
 M(z):=z(3+z^2)I-L-3z^2H.
\]
The coefficient of $z^{3K}$ in $\det M(z)$ is $1$: it is obtained uniquely
by selecting the $z^3$ term from every diagonal entry.  Thus $\det M$ is not
the zero polynomial.  For all but finitely many $z\in(0,1)$, $M(z)$ is
invertible and therefore $G(z)=0$, contradicting positivity of probability
PGFs on $(0,1)$.
\end{proof}

\begin{theorem}[No positive scalar state-only drift]
\label{thm:noscalardrift}
Suppose $\phi:\Z_{\ge0}\to\R$ and constants $c,\lambda_x,\mu_y$ satisfy
\begin{equation}
 \phi(j+y-x)-\phi(j)
 \ge c+\lambda_x+\mu_y
 \label{eq:scalardrift}
\end{equation}
whenever $j+y-x\ge0$, with gauges
\[
 \sum_xP(x)\lambda_x=0,
 \qquad
 \sum_yQ(y)\mu_y=0.
\]
For the dyadic pair, $c\le0$.
\end{theorem}

\begin{proof}
For the dyadic pair the difference \(d:=y-x\) takes values \(-1,1,3\).
The \(+1\) inequalities give a uniform lower bound on
\(d_j:=\phi(j+1)-\phi(j)\), while the \(-1\) inequality gives a uniform upper
bound.  Hence \((d_j)\) is bounded.

Set \(a_N=N^{-1}\sum_{j=1}^N d_j\).  Along a subsequence
\(N_k\to\infty\), \(a_{N_k}\to a\).  For each fixed
\(r\in\{-1,1,3\}\), boundedness implies
\begin{equation}
 \frac1{N_k}\sum_{j=1}^{N_k}[\phi(j+r)-\phi(j)]\longrightarrow ra.
 \label{eq:cesaro-shift}
\end{equation}
For positive \(r\), expand into consecutive first differences; for \(r=-1\),
shift the index.  The finitely many endpoint terms are negligible because
\((d_j)\) is bounded.

Average \eqref{eq:scalardrift} over \(j=1,\ldots,N_k\) and over
\((X,Y)\sim P\otimes Q\).  All moves are admissible for \(j\ge1\), and the
gauge terms vanish.  Passing to the subsequential limit using
\eqref{eq:cesaro-shift} gives
\[
 a\,\E(Y-X)\ge c.
\]
The two increment laws have equal mean, so \(\E(Y-X)=0\), and therefore
\(c\le0\).
\end{proof}

These theorems have a useful positive interpretation.  A sharp asymptotic
object, if it exists, must retain information at growing scales.  The large
kernel $W_m$ is not an artifact of a poor coordinate system that can be
compressed into finitely many exact modes without loss.

% ============================================================
\section{The renormalized positive phase problem}
\label{sec:renormproblem}
% ============================================================

We can now state the unresolved problem without appealing to a particular
buffer ansatz.

The linear quotient $V_m=U_m\oplus W_m$ separates visible total-information
profiles from ordering information.  Theorem~\ref{thm:positivefiber} describes
the complete positive fiber before causality, while
Theorem~\ref{thm:causal-farkas} gives a complete finite-dimensional dual
criterion for the local causal matching fiber at every fixed block length.
Thus the unresolved issue is now more specific: \emph{future zero-repair
viability must be incorporated in a way that is stable under scale}.

The role of Section~\ref{sec:doublingphase} is therefore decisive.  Blocking from $m$
to $2m$ does not merely transport an old phase vector.  It creates the fresh
sector $K_m$ of dimension $m^2$.  Thus any exact scale map must account for at
least three operations:
\begin{center}
\fbox{\parbox{0.78\textwidth}{\centering
transport inherited phase $+$ create inter-block phase\\[1mm]
$+$ project back into the positive causal fiber.}}
\end{center}
The last operation is nonlinear because positivity and future viability are
inequalities.

\subsection{The central characterization problem}

At fixed \(m\), the local coupling fiber is no longer mysterious:
Theorem~\ref{thm:causal-farkas} characterizes it exactly, and a finite set of
rational dual rays exists in principle.  What is missing is the analogue of a
single structural criterion whose form survives scale doubling and already
incorporates descendant viability.

\begin{problem}[Scale-stable zero-repair phase characterization]
\label{prob:phasecharacterization}
Find coarse/phase coordinates and a family of necessary-and-sufficient
positivity or dual inequalities such that: (i) one-scale lifts land in child
states viable for the prescribed remaining horizon; (ii) the form closes under
\(m\mapsto2m\); and (iii) the induced repair cost is uniformly bounded in the
dyadic scale.  Equivalently, compress the exponentially expanding fixed-block
Farkas descriptions into an intrinsic renormalization law on the positive
zero-repair phase state.
\end{problem}

For \(m=2\), Proposition~\ref{prop:lift} isolates the unique one-sided phase
coordinate \(\Theta\).  For general \(m\),
Theorem~\ref{thm:doublingdecomp} shows why a single \(\Theta\) cannot suffice,
and Theorem~\ref{thm:causal-farkas} shows that the obstacle is not
finite-dimensional duality itself but the absence of a scale-uniform
compression.

\subsection{A conditional scale theorem}

The following simple implication isolates exactly what an upper-bound theorem
would have to deliver.

\begin{proposition}[Bounded scale repair implies logarithmic growth]
\label{prop:boundedscale}
Suppose there is a constant $C<\infty$ such that
\begin{equation}
 \boxed{
 \cB_{2n}\le\cB_n+C
 }
 \label{eq:doublingupper}
\end{equation}
for every $n\ge1$.  Then
\begin{equation}
 \cB_n=O(\log n).
 \label{eq:logupper}
\end{equation}
For the dyadic parity pair, this combines with
Theorem~\ref{cor:one-third} to give
\begin{equation}
 \boxed{\cB_n=\Theta(\log n).}
 \label{eq:theta-log-conditional}
\end{equation}
\end{proposition}

\begin{proof}
Iterating \eqref{eq:doublingupper} gives
$\cB_{2^k}\le\cB_1+kC$.  Monotonicity of the finite-horizon viability bill
then gives the same order for arbitrary $n$ by choosing
$2^{k-1}<n\le2^k$.  The lower bound is Theorem~\ref{cor:one-third}.
\end{proof}

Equation~\eqref{eq:doublingupper} is not proved here.  It is the most concrete
sharp upper target suggested by the scale algebra.  If the dyadic increment
converged,
\[
 \cB_{2n}-\cB_n\longrightarrow\gamma,
\]
then along dyadic horizons one would obtain
\[
 \cB_n\sim\frac{\gamma}{\log2}\log n,
\]
and the rigorous lower bound would force
\begin{equation}
 \gamma\ge\frac{\log2}{3}.
 \label{eq:gammalower}
\end{equation}
No equality is asserted.

\subsection{Why the one-channel Jordan picture is insufficient}

The explicit packet $\eta_s$ is marginal in the second factorial moment, but
its first-moment carrying cost decays geometrically: an envelope mass of order
$4^{-s}$ transported to distance $2^s$ costs order $2^{-s}$.  This is exactly
why Theorem~\ref{thm:envelope} has finite mean.  Consequently, one copy of the
explicit $\eta$ channel cannot by itself explain logarithmic growth.

A more plausible critical mechanism, if the logarithmic lower bound is sharp,
is the repeated creation of fresh phase directions under blocking.  The
ambient fresh sector has dimension $m^2$, while the cubic signed residual is
suppressed by a third-order scale factor.  The dimensional balance
\begin{equation}
 m^2\cdot m^{-3}\cdot m\asymp1
 \label{eq:dimensionalbalance}
\end{equation}
is suggestive: multiplicity, cubic residual mass, and first-moment transport
distance cancel at order three.  We emphasize that
\eqref{eq:dimensionalbalance} is a scaling heuristic, not a theorem.  It does
not show that $\Theta(m^2)$ feasible phase directions are independently
excited, nor that positivity prevents cancellations.  Its value is that it
identifies a falsifiable mechanism compatible with both the exact cubic symbol
and the exact growth of the phase kernel.

% ============================================================
\section{Relation to neighboring mathematical theories}
\label{sec:neighbors}
% ============================================================

\paragraph{Adapted and causal optimal transport.}
The causal constraint in Section~\ref{sec:5over16} is a standard
nonanticipativity constraint on a coupling.  General causal and bicausal
transport duality, including dual attainment under broad measurable-cost
hypotheses, is available in \cite{KrsekPammer2025}.  Very recent graph-causal
optimal transport characterizes gluing properties and dynamic programming for
causality prescribed by directed acyclic graphs \cite{OblojTuchilus2026}.
These results make adapted transport a natural surrounding language.  They do
not make the present problem automatic: exact inventory conservation and the
hard cone constraint $J\ge0$ must still be incorporated, and the full
countable-state no-gap theorem for our Bellman--Fenchel recursion has not been
proved.

\paragraph{Convex dynamic programming.}
The passage from convex viability sets to lower support functionals is a
finite-dimensional convex duality calculation, not a new general duality
theorem.  For infinite-dimensional stochastic optimization with information
constraints, the recent work of Pennanen and Perkki\"o provides a broader
convex-duality and dynamic-programming framework \cite{PennanenPerkkio2026Dual,PennanenPerkkio2026DP}.
The unresolved issue here is to choose a weighted function/measure topology in
which the unbounded mean objective, shift invariance and positivity coexist
with strong duality.

\paragraph{Nonlinear Perron--Frobenius theory.}
Order-preserving homogeneous and additively homogeneous maps, quotienting by
constants, and projective metrics are standard in nonlinear
Perron--Frobenius theory and in Shapley operators
\cite{LemmensNussbaum2012,Lins2023,AkianGaubertHochart2020}.  Recent work also
extends joint spectral-radius ideas to switched nonlinear order-preserving
systems \cite{DeiddaGuglielmiTudisco2025}.  The relevance here is structural,
not a theorem transfer: Theorem~\ref{thm:nofiniteautomaton} already rules out
an exact finite-mode state.  Across all scales the zero-repair phase state is
necessarily infinite-dimensional and constrained by positivity.

\paragraph{Subdivision and vanishing moments.}
The identity
\[
 \Delta_3\cR=\frac14D_2\Delta_3
\]
places the boundary algebra directly next to subdivision schemes.  The mask
$(1+z)^3/4$ is the binary quadratic B-spline mask; its order-three zero
condition is the same algebra that annihilates the matched mass, mean and
variance jets.  Sum rules, polynomial reproduction and invariant difference
spaces are classical tools in subdivision theory \cite{CharinaConti2013}.
The novelty of the present difficulty is the additional causal-positive fiber
sitting over this linear scheme.

\paragraph{Finitary coding and randomness recycling.}
The zero-repair equation resembles a constrained information coboundary, while
finitary isomorphism theory shows that stronger moment assumptions on coding
length preserve increasingly fine information invariants
\cite{Parry1979,Schmidt1984,HarveyPeres2011}.  In a different online-sampling
model, modern randomness-recycling algorithms can approach the Shannon limit
using persistent state; recent real-probability algorithms achieve expected
entropy $\sum_iH(F_i)+O(\log n)$ with $O(\log n)$ persistent space
\cite{DraperSaad2026,DraperHarrisSaad2026}.  Neither result implies a
zero-repair theorem, but both reinforce the need to distinguish entropy rate
from persistent working state.

\paragraph{Synchronization-channel duality.}
Kova\v{c}evi\'c's recent exact sticky-channel calculation is a useful nearby
example of proof architecture rather than a source of a theorem used here:
a channel is reduced to a simpler operational object, a dual bound is made
exact by a coefficientwise criterion, and the entire domination class is then
characterized algebraically \cite{Kovacevic2026}.  Problem
\ref{prob:phasecharacterization} asks for the analogous structural step for a
causal-positive phase fiber.

% ============================================================
\section{Back to ``Who pays the bill?''}
\label{sec:original}
% ============================================================

The original resource question was broader than the zero-repair model.  An
observable stochastic system may be realizable at the ordinary Shannon rate
while different exact-realization contracts impose different requirements on
its hidden state.  The mathematics above separates three quantities that
should not be conflated:
\begin{equation}
 \boxed{
 \text{entropy rate}
 \quad\neq\quad
 \text{fresh repair expenditure}
 \quad\neq\quad
 \text{exact causal working inventory}.
 }
 \label{eq:threebills}
\end{equation}

The first is a rate: how many new random bits per step are required
asymptotically.  The second measures information that a simulator is permitted
to inject after seeing where mismatch occurred.  The third is the present
quantity: a nonnegative inventory that must make every one-step information
identity exact without such fresh repair.

For the dyadic parity pair, the present paper proves a sharp qualitative
statement about the third resource.  Even though mass, mean and variance of
the two information laws agree, their third-order mismatch cannot be carried
forever by any finite-mean zero-repair inventory:
\begin{equation}
 \cB_n\ge\left(\frac13-o(1)\right)\log n.
 \label{eq:backbill}
\end{equation}
At the same time, the two-symbol causal matching theorem shows why a scalar
information spectrum alone is insufficient: causality creates an additional
positive obstruction while leaving the signed scalar residual unchanged.
Thus part of the bill is carried by ordering information that disappears under
the usual scalar projection.

This does \emph{not} prove that an unrestricted causal simulator pays an
additional asymptotic entropy rate.  In the broader causal-tree model, fresh
randomness, nonlocal blocks and hidden dependence may implement strategies
forbidden by zero repair.  A companion analysis of that broader contract finds
that, for two finite actions, the first-order causal rate can reach the
observable Shannon floor \cite{MikulikCausal}.  There is no contradiction:
rate and working inventory are different resources.

The missing bridge is therefore mathematically precise.

\begin{problem}[Entropy--inventory comparison]
\label{prob:entropyinventory}
Establish the sharpest general comparison, if any, between the finite-horizon
unrestricted causal entropy residual and the zero-repair viability bill
$\cB_n$.  In particular, determine which forms of fresh repair or blockwise
hidden dependence can convert a diverging working-inventory requirement into
only $o(n)$ additional source entropy.
\end{problem}

The answer would return the phase-renormalization theory developed here to the
original accounting problem.  Until such a bridge is proved, the correct
conclusion is operational rather than ontological: exact realization costs
depend on what must remain simultaneously valid, on what may be repaired
later, and on which history information the hidden state is allowed to retain.

% ============================================================
\section{Theorem ledger and open frontier}
\label{sec:ledger}
% ============================================================

For clarity, Table~\ref{tab:ledger} separates the proved statements from the
renormalization program.

\begin{table}[ht]
\centering
\small
\begin{tabular}{p{0.72\textwidth}p{0.18\textwidth}}
\toprule
Statement & Status\\
\midrule
Finite-horizon fixed-tilt logarithmic converse & proved\\
Third-cumulant lower bound $\liminf\cB_n/\log n\ge\Delta\kappa_3/(6\sigma^2)$ & proved\\
Dyadic lower coefficient $1/3$ and $\cB_n\to\infty$ & proved\\
Finite-dimensional Bellman--Fenchel duality & proved\\
Countable-state Bellman formula as weak dual certificate & proved\\
Phase quotient $W_m=\ker\Pi_m$, $\dim W_m=2^m-(m+1)$ & proved\\
Exact positive scalar-profile fiber, Theorem~\ref{thm:positivefiber} & proved\\
Fixed-block local causal phase fiber: exact polyhedral/Farkas characterization & proved\\
Two-symbol projected causal profile: exact two-facet characterization & proved\\
Two-symbol causal mismatch $5/16$ with solver-free proof and exact primal--dual certificate & proved\\
Causal surcharge invisible to the signed scalar residual & proved\\
Block-doubling decomposition of $W_{2m}$ and fresh sector $\dim K_m=m^2$ & proved\\
Cubic shift factorization and Green affine conjugacy & proved\\
Explicit marginal phase cohomology and mean-$17/6$ envelope & proved$^*$\\
Mahler scale fixed point with finite first and infinite second moment & proved\\
Finite exact profile automaton & impossible\\
Positive scalar state-only Bellman drift & impossible\\
Scale-stable characterization of the zero-repair causal phase state & open\\
Strong countable-state Bellman--Fenchel duality in a natural weighted space & open\\
Uniform doubling bound $\sup_n(\cB_{2n}-\cB_n)<\infty$ & open\\
Matching $O(\log n)$ upper bound / $\Theta(\log n)$ law & open\\
Exact asymptotic constant for $\cB_n/\log n$ & open\\
General comparison with unrestricted causal entropy residual & open\\
\bottomrule
\end{tabular}
\caption{Logical status of the main statements. The symbol $^*$ means: proved for the explicit phase channel stated in Section~\ref{sec:eta}, not for the full zero-repair phase dynamics.}
\label{tab:ledger}
\end{table}

The frontier can be summarized in one sentence:
\begin{center}
\fbox{\parbox{0.78\textwidth}{\centering
the cubic linear defect and every fixed-block causal fiber are explicit;\\[1mm]
the remaining object is their scale-stable zero-repair renormalization.}}
\end{center}

% ============================================================
\section{Conclusion}
% ============================================================

The first nontrivial dyadic zero-repair problem contains two distinct kinds of
information.  The scalar information profile records how much surprisal has
accumulated; the phase fiber records in which order that information was
created.  Static convolution sees only the first.  Causality sees both.

This distinction is already completely visible after two symbols.  The
projection of the causal coupling polytope is the static scalar transport
polytope cut by two explicit oblique facets.  One of these facets alone forces
the static mismatch $1/4$ up to the exact causal value $5/16$.  The additional
$1/16$ appears identically on both scalar residuals and is therefore invisible
to their signed difference.  At larger block lengths the hidden space expands
rapidly: $W_m$ has dimension $2^m-(m+1)$, and every doubling creates a fresh
$m^2$-dimensional anti-diagonal sector even before positivity is imposed.

Independently, the source mismatch has an exact cubic normal form.  The factor
$(I-E)^3$ annihilates the matched mass, mean and variance jets; division by the
corresponding cubic zero produces a Green coordinate with a positive unit
source and a quadratic B-spline boundary mask.  The same order three is
critical in the rigorous Laplace converse, which forces
\[
 \cB_n\ge\left(\frac13-o(1)\right)\log n
\]
for the dyadic pair.

The two structures meet at the unresolved point.  Linear cubic inversion and
fixed-block causal feasibility are explicit; keeping all descendants inside a
positive, prefix-compatible \emph{future-viable} phase state is not.  A
scale-stable characterization of that state would turn the present
renormalization picture into an actual upper-bound mechanism.  A bounded cost per dyadic scale would close the growth
order at $\Theta(\log n)$.

The broader information-theoretic lesson is narrower but already rigorous:
an exact stochastic realization can have different costs depending on whether
one prices source entropy, fresh repair, or nonnegative causal working
inventory.  The scalar information spectrum does not contain all of the
causal bookkeeping.  Ordering information can carry a bill of its own.

% ============================================================
\appendix
\section{Complete rational certificate for the two-symbol causal LP}
\label{app:5over16matrix}
% ============================================================

This appendix gives an independent full rational LP certificate for Corollary~\ref{thm:5over16}.  The closed-form two-facet proof in Section~\ref{sec:5over16} does not require this appendix, but the certificate makes the original sixteen-variable formulation fully checkable without a solver.
Vectorize the coupling matrix row by row as
\[
 \pi=(\pi_{00,00},\pi_{00,01},\ldots,\pi_{11,11})^T.
\]
With the constraint order used in the proof, the equality matrix is
\begin{equation}
\resizebox{\textwidth}{!}{$
A=\begin{pmatrix}
1&1&1&1&0&0&0&0&0&0&0&0&0&0&0&0\\
0&0&0&0&1&1&1&1&0&0&0&0&0&0&0&0\\
0&0&0&0&0&0&0&0&1&1&1&1&0&0&0&0\\
0&0&0&0&0&0&0&0&0&0&0&0&1&1&1&1\\
1&0&0&0&1&0&0&0&1&0&0&0&1&0&0&0\\
0&1&0&0&0&1&0&0&0&1&0&0&0&1&0&0\\
0&0&1&0&0&0&1&0&0&0&1&0&0&0&1&0\\
0&0&0&1&0&0&0&1&0&0&0&1&0&0&0&1\\
\frac14&-\frac34&0&0&\frac14&-\frac34&0&0&0&0&0&0&0&0&0&0\\
-\frac14&\frac34&0&0&-\frac14&\frac34&0&0&0&0&0&0&0&0&0&0\\
0&0&\frac14&-\frac34&0&0&\frac14&-\frac34&0&0&0&0&0&0&0&0\\
0&0&-\frac14&\frac34&0&0&-\frac14&\frac34&0&0&0&0&0&0&0&0\\
0&0&0&0&0&0&0&0&\frac14&-\frac34&0&0&\frac14&-\frac34&0&0\\
0&0&0&0&0&0&0&0&-\frac14&\frac34&0&0&-\frac14&\frac34&0&0\\
0&0&0&0&0&0&0&0&0&0&\frac14&-\frac34&0&0&\frac14&-\frac34\\
0&0&0&0&0&0&0&0&0&0&-\frac14&\frac34&0&0&-\frac14&\frac34
\end{pmatrix}.$}
\label{eq:fullA}
\end{equation}
The right-hand side is \eqref{eq:5over16b}.  The first eight rows impose the
two marginals; the last eight are the causal equations
\eqref{eq:causalLPconstraints}.  Redundant rows are retained to match the
lexicographic construction.

For the dual vector \eqref{eq:dualvector}, direct multiplication gives
\[
 A^Ty=\begin{pmatrix}
1&1&1&-1\\0&0&0&-2\\0&-1&-1&0\\1&0&0&1
\end{pmatrix},
\qquad
 C=\begin{pmatrix}
1&1&1&1\\0&1&1&1\\0&1&1&1\\1&0&0&1
\end{pmatrix}.
\]
Thus \(A^Ty\le C\) componentwise.  Finally,
\[
 b^Ty=\frac1{16}\left(1-\frac34+\frac{27}{4}-2\right)=\frac5{16}.
\]
Together with the primal matrix \eqref{eq:primalmatrix}, this is an exact
rational primal--dual certificate of the optimum.

% ============================================================
\section*{Research-status and reproducibility note}
% ============================================================

All asymptotic claims in this paper are supported by analytic arguments.
The exact two-symbol causal value is accompanied by explicit rational primal
and dual certificates, so it does not depend on floating-point optimization.
Computer algebra was used to check finite polynomial identities and exact
coefficients during development, not as a substitute for the stated proofs.

The paper deliberately uses conservative priority language.  The surrounding
subjects---adapted transport, finitary coding, nonlinear positive operators,
subdivision theory and synchronization channels---are mature, and an
equivalent formulation of some structural lemma may exist elsewhere.  Where a
direct precedent is known, it is cited; no stronger historical novelty claim
is intended.

During preparation, the author used OpenAI's ChatGPT as a research assistant
for literature discovery, algebraic checking, adversarial proof review and
\LaTeX{} preparation.  The mathematical statements are separated explicitly
into proved results and open problems in Table~\ref{tab:ledger}.

% ============================================================
\begin{thebibliography}{99}

\bibitem{Rockafellar1970}
R.~T.~Rockafellar,
\emph{Convex Analysis},
Princeton University Press, 1970.

\bibitem{Parry1979}
W.~Parry,
\emph{Finitary isomorphisms with finite expected code lengths},
Bull. London Math. Soc. \textbf{11} (1979), 170--176.
\url{https://doi.org/10.1112/blms/11.2.170}

\bibitem{Schmidt1984}
K.~Schmidt,
\emph{Invariants for finitary isomorphisms with finite expected code lengths},
Invent. Math. \textbf{76} (1984), 33--40.

\bibitem{HarveyPeres2011}
N.~Harvey and Y.~Peres,
\emph{An invariant of finitary codes with finite expected square root coding length},
Ergodic Theory Dynam. Systems \textbf{31} (2011), 77--90.
\url{https://doi.org/10.1017/S014338570900090X}

\bibitem{LemmensNussbaum2012}
B.~Lemmens and R.~D.~Nussbaum,
\emph{Nonlinear Perron--Frobenius Theory},
Cambridge Tracts in Mathematics, Vol.~189, Cambridge University Press, 2012.

\bibitem{CharinaConti2013}
M.~Charina and C.~Conti,
\emph{Polynomial reproduction of multivariate scalar subdivision schemes},
J. Comput. Appl. Math. \textbf{240} (2013), 51--61.
\url{https://doi.org/10.1016/j.cam.2012.06.013}

\bibitem{AkianGaubertHochart2020}
M.~Akian, S.~Gaubert, and A.~Hochart,
\emph{A game theory approach to the existence and uniqueness of nonlinear
Perron--Frobenius eigenvectors},
Discrete Contin. Dyn. Syst. \textbf{40} (2020), 207--231.
\url{https://doi.org/10.3934/dcds.2020009}

\bibitem{Lins2023}
B.~Lins,
\emph{A unified approach to nonlinear Perron--Frobenius theory},
Linear Algebra Appl. \textbf{675} (2023), 48--89.
\url{https://doi.org/10.1016/j.laa.2023.06.018}

\bibitem{KrsekPammer2025}
D.~Kr\v{s}ek and G.~Pammer,
\emph{General duality and dual attainment for adapted transport},
Appl. Math. Optim. \textbf{91} (2025), Article 52.
\url{https://doi.org/10.1007/s00245-025-10240-y}

\bibitem{DeiddaGuglielmiTudisco2025}
P.~Deidda, N.~Guglielmi, and F.~Tudisco,
\emph{Nonlinear joint spectral radius},
arXiv:2507.11314, 2025.

\bibitem{PennanenPerkkio2026Dual}
T.~Pennanen and A.-P.~Perkki\"o,
\emph{Duality in convex stochastic optimization},
Math. Program. \textbf{215} (2026), 145--191.
\url{https://doi.org/10.1007/s10107-025-02216-1}

\bibitem{PennanenPerkkio2026DP}
T.~Pennanen and A.-P.~Perkki\"o,
\emph{Dynamic programming and dimensionality in convex stochastic optimization
and control},
Optim. Lett. \textbf{20} (2026), 257--273.
\url{https://doi.org/10.1007/s11590-025-02230-4}

\bibitem{DraperSaad2026}
T.~L.~Draper and F.~A.~Saad,
\emph{Efficient online random sampling via randomness recycling},
in Proc. 2026 ACM--SIAM Symposium on Discrete Algorithms (SODA),
2026, pp.~2473--2511.
\url{https://doi.org/10.1137/1.9781611978971.89}

\bibitem{DraperHarrisSaad2026}
T.~L.~Draper, D.~G.~Harris, and F.~A.~Saad,
\emph{Online random sampling with real probabilities},
arXiv:2607.13828, 2026.

\bibitem{Kovacevic2026}
M.~Kova\v{c}evi\'c,
\emph{The capacity of a family of sticky channels},
arXiv:2607.28281, 2026.

\bibitem{OblojTuchilus2026}
J.~Ob\l\'oj and V.~Tuchilus,
\emph{Graph causal optimal transport and Wasserstein distances},
arXiv:2608.13716, 2026.

\bibitem{MikulikCausal}
J.~Mikul\'ik,
\emph{Who Pays the Bill? Counterfactual Identity, the Causal Shannon Floor,
Spectral Ledgers, and Orbit Resolution},
research preprint, 2026.

\end{thebibliography}

\end{document}
